\documentclass{article}


% if you need to pass options to natbib, use, e.g.:
%     \PassOptionsToPackage{numbers, compress}{natbib}
% before loading neurips_2025


% ready for submission
% \usepackage{neurips_2025}


% to compile a preprint version, e.g., for submission to arXiv, add add the
% [preprint] option:
% \usepackage[preprint]{neurips_2025}


% to compile a camera-ready version, add the [final] option, e.g.:
%     \usepackage[final]{neurips_2025}


% to avoid loading the natbib package, add option nonatbib:
\usepackage[nonatbib,preprint]{neurips_2025}



\usepackage[utf8]{inputenc} % allow utf-8 input
\usepackage[T1]{fontenc}    % use 8-bit T1 fonts
\usepackage[colorlinks,
            linkcolor=red,
            anchorcolor=blue,
            citecolor=cyan]{hyperref} 
\usepackage{url}            % simple URL typesetting
\usepackage{booktabs}       % professional-quality tables
\usepackage{amsfonts}       % blackboard math symbols
\usepackage{nicefrac}       % compact symbols for 1/2, etc.
\usepackage{microtype}      % microtypography
% \usepackage{xcolor}         % colors
\usepackage[table,xcdraw]{xcolor}
\usepackage{adjustbox}
\usepackage{subfigure}
\usepackage{multirow,multicol}
\usepackage{arydshln}
\usepackage{wrapfig}    % colors
\usepackage{fontawesome}
\usepackage{pifont}
\usepackage{amsmath}

\definecolor{lightgray}{gray}{0.9}
\newcommand{\CLA}[1]{{\color[HTML]{4472c4} \textbf{#1}}}
\newcommand{\CLB}[1]{{\color[HTML]{E76254} \textbf{#1}}}

\newcommand{\GroupA}[1]{{\color[HTML]{92ca2c} \textbf{#1}}}
\newcommand{\GroupB}[1]{{\color[HTML]{3c7daa} \textbf{#1}}}
\newcommand{\GroupC}[1]{{\color[HTML]{e64f04} \textbf{#1}}}
\newcommand{\GroupD}[1]{{\color[HTML]{f2c400} \textbf{#1}}}
\newcommand{\GroupE}[1]{{\color[HTML]{a757a7} \textbf{#1}}}
\newcommand{\GroupF}[1]{{\color[HTML]{56B1AD} \textbf{#1}}}
\newcommand{\GroupG}[1]{{\color[HTML]{222222} \textbf{#1}}}

\newcommand{\mat}[1]{\begin{pmatrix} #1 \end{pmatrix}}
\newcommand{\T}{^\mathsf{T}}
\newcommand{\inv}{^{-1}}
\newcommand{\RotX}[1]{\mathbf{R}_x(#1)}
\newcommand{\RotY}[1]{\mathbf{R}_y(#1)}
\newcommand{\RotZ}[1]{\mathbf{R}_z(#1)}
\newcommand{\TransX}[1]{\mathbf{Tr}_x(#1)}
\newcommand{\TransY}[1]{\mathbf{Tr}_y(#1)}
\newcommand{\TransZ}[1]{\mathbf{Tr}_z(#1)}
\newcommand{\s}[1]{\sin(#1)}
\newcommand{\cs}[1]{\cos(#1)}
\newcommand{\sa}[1]{\sin(\alpha_{#1})} % \alpha_{i-1} typically
\newcommand{\ca}[1]{\cos(\alpha_{#1})} % \alpha_{i-1} typically
\newcommand{\st}[1]{\sin(\theta_{#1})}
\newcommand{\ct}[1]{\cos(\theta_{#1})}
\newcommand{\stt}[2]{\sin(\theta_{#1}+\theta_{#2})}
\newcommand{\ctt}[2]{\cos(\theta_{#1}+\theta_{#2})}
\newcommand{\atan}[2]{\operatorname{atan2}(#1, #2)}

\title{Image Quality Assessment for Embodied AI}


% The \author macro works with any number of authors. There are two commands
% used to separate the names and addresses of multiple authors: \And and \AND.
%
% Using \And between authors leaves it to LaTeX to determine where to break the
% lines. Using \AND forces a line break at that point. So, if LaTeX puts 3 of 4
% authors names on the first line, and the last on the second line, try using
% \AND instead of \And before the third author name.


\author{\hspace{-7pt} \textbf{Chunyi Li}$^{1,2,3}$ \textbf{Jiahao Xiao}$^{1,2}$  \textbf{Jianbo Zhang}$^{1}$  \textbf{Farong Wen}$^{1,2}$  \textbf{Zicheng Zhang}$^{1,2}$ \textbf{Yuan Tian}$^{1,2}$ \\ \textbf{Xiangyang Zhu}$^{2}$ \hspace{1pt} \textbf{Xiaohong Liu}$^{1}$ \hspace{1pt} \textbf{Zhengxue Cheng}$^{1}$\hspace{1pt} \textbf{Weisi Lin}$^{3}$ \hspace{1pt} \textbf{Guangtao Zhai}$^{1,2}$ \vspace{2pt} \\
$^1$ Shanghai Jiao Tong University \hspace{9pt} $^2$ Shanghai AI Lab  \hspace{9pt} $^3$ Nanyang Technological University \vspace{2pt}\\
Project Page: \href{https://github.com/lcysyzxdxc/EmbodiedIQA}{\textit{https://github.com/lcysyzxdxc/EmbodiedIQA}}
}
  % examples of more authors
  % \And
  % Coauthor \\
  % Affiliation \\
  % Address \\
  % \texttt{email} \\
  % \AND
  % Coauthor \\
  % Affiliation \\
  % Address \\
  % \texttt{email} \\
  % \And
  % Coauthor \\
  % Affiliation \\
  % Address \\
  % \texttt{email} \\
  % \And
  % Coauthor \\
  % Affiliation \\
  % Address \\
  % \texttt{email} \\


\begin{document}


\maketitle

\begin{figure}[tbph]
    \centering
    \vspace{-15pt}
    \includegraphics[width=\linewidth]{figs/spotlight.pdf}
    \vspace{-6pt}
    \caption{The significant gap between human, machine, and robot visual systems. Humans and Machines are sensitive to different distortions, while Robots have \CLA{Decision} and \CLA{Execution} steps beyond \CLA{Cognition}, highlighting the importance of a \CLA{Perception} quality index for Embodied AI.}
    \label{fig:spotlight}
\end{figure}

\begin{abstract}
Embodied AI has developed rapidly in recent years, but it is still mainly deployed in laboratories, with various distortions in the Real-world limiting its application. Traditionally, Image Quality Assessment (IQA) methods are applied to predict human preferences for distorted images; however, there is no IQA method to assess the usability of an image in embodied tasks, namely, the perceptual quality for robots. To provide accurate and reliable quality indicators for future embodied scenarios, we first propose the topic: IQA for Embodied AI. Specifically, we (1) based on the Mertonian system and meta-cognitive theory, constructed a perception-cognition-decision-execution pipeline and defined a comprehensive subjective score collection process; (2) established the Embodied-IQA database, containing over 30k reference/distorted image pairs, with more than 5m fine-grained annotations provided by Vision Language Models/Vision Language Action-models/Real-world robots; (3) trained and validated the performance of mainstream IQA methods on Embodied-IQA, demonstrating the need to develop more accurate quality indicators for Embodied AI. We sincerely hope that through evaluation, we can promote the application of Embodied AI under complex distortions in the Real-world.
\end{abstract}


\section{Introduction}

To achieve Artificial General Intelligence (AGI), Embodied AI, as a bridge connecting external and internal realities, has developed rapidly in recent years. Relying on its ability to interact with the physical environment, Embodied AI \cite{review:embodied1,review:embodied2,review:embodied3} has been applied to simple scenarios such as factories, warehouses, and household chores, but it is not yet capable of handling complex environments like autonomous driving and wilderness exploration. Unlike traditional robotics driven by fixed algorithms, Embodied AI collects signals from the Real-world and is therefore susceptible to distortions. For example, a pick-and-place task may be successfully debugged in the laboratory, but it may fail in Real-world applications due to slight lens defocusing or shaking. Therefore, the preferences of Embodied AI should be analyzed to filter out these low-quality images.

For human viewers, this problem can be solved through Image Quality Assessment (IQA) metrics. For example, in streaming media, collect human subjective preferences for distorted images. Since human resources are expensive, IQA will develop objective quality indicators to fit subjective scores. Similarly, for Embodied AI, it is also necessary to collect the success rate of downstream tasks using distorted images, quantifying the fidelity with the reference results, and developing IQA metrics.

Unfortunately, due to the significant differences between Human/Machine/Robot Visual Systems (HVS/MVS/RVS), previous IQA methods cannot be directly transferred to Embodied AI scenarios, as shown in Figure \ref{fig:spotlight}. First, HVS and MVS are sensitive to different types of distortions. Humans are sensitive to distortions such as noise and compression, which do not affect the downstream tasks of machines. Brightness and contrast are the opposite. Therefore, as a machine, Embodied AI cannot use the past human-oriented processing methods. Second, although MVS and RVS are sensitive to some common distortions, the perceptual quality of a general machine only depends on the performance of segmentation and detection tasks that belong to Cognition. Robots, however, have subsequent Decision and Execution steps. High fidelity in the previous step does not guarantee the success of the next. Therefore, unlike a general machine, it is necessary to fully consider Cognition, Decision, and Execution to characterize the Perception of Embodied AI. Considering these issues, we first attempt to implement IQA metrics into Embodied AI. Our contributions are summarized as follows:
\begin{itemize}
    \vspace{-2pt}
    % \item Theory: We refer to Merton Law in robotic intelligence to construct the Perception-Cognition-Decision-Execution pipeline. We define tasks related to Embodied Perception and specify the subjects for each step in Cognition, Decision, and Execution, as well as the process for collecting subjective scores. 
    \item Theory: We refer to the Mertonian Law in robotic intelligence to construct the Perception-Cognition-Decision-Execution pipeline. We define tasks related to Embodied Perception and specify the subjects for each step in Cognition, Decision, and Execution. 
    \item Data: We add corruption to images in Embodied tasks, collecting over 36k reference/distorted image pairs. We perform inference using Vision Language Models (VLM) and Vision Language Action-model (VLA) for over 5 million annotations. This large-scale database can effectively drive the development of quality metrics for Embodied AI.
    \item Experiment: We experiment with 15 advanced IQA methods on our database, proving that more sophisticated IQA metrics are needed for Embodied AI. Additionally, we first conduct real-world experiments in the IQA field, executing 1.5k Embodied tasks in the Real-world, revealing the internal connections between Cognition, Decision, and Execution.
    \vspace{-4pt}
\end{itemize}
 

\section{Related Works}
\vspace{-4pt}
\subsection{Mertonian System for Robotic Intelligence}

Intelligent models are divided into Newtonian and Mertonian \cite{review:merton} systems. Newtonian systems typically refer to those systems that can be precisely described and predicted by deterministic physical laws, such as classical mechanical systems. In contrast, Mertonian systems involve systems that include `free will', whose behavior is influenced by feedback between beliefs and actions. The characteristic of such systems is that even given the current state and control conditions, the next state of the system cannot be accurately obtained by solving, so its behavior is difficult to predict precisely.

HVS and MVS can both be simplified as Newtonian systems, since their Decision and Execution processes are robust. After humans solve a problem, the body can accurately execute the will; machines like computers will also output results on the screen after completing segmentation/detection algorithms. However, the Decision and Execution of RVS do not fully match Cognition. For example, a deviation of one character in Cognition may greatly change the pose in Decision; a one-centimeter path offset in Decision may also cause Execution to hit obstacles. Since the impact of distortion on them is unpredictable, it is necessary to handle these three steps separately for the IQA task.

\subsection{Image Quality Assessment for Machine}

Since 1999, Perception has been recognized as the first step in the interaction between AI agents and external reality, whose mechanism \cite{review:perception1,review:perception2,review:perception3,review:perception4} has been revealed. However, no perceptual quality score has been assigned to each distorted image like current IQA \cite{dataset:cmc,my:rbench,add:aigc} metrics, which is exactly Embodied AI needs in Real-world applications. In the past decades, IQA has been widely studied as shown in Table \ref{tab:database}, but none of them meet the above needs of Embodied AI. First, most databases are human-oriented, with only two coarse-grained \cite{add:smr1,add:smr2}, two fine-grained \cite{dataset:mpd,dataset:epd} machines as subjects (VLM and reinforcement learning); second, as mentioned above, no database covers Cognition, Decision, and Execution altogether. Considering the characteristics of Embodied AI, it is necessary to establish a new dataset in accordance with the requirements of the Mertonian system.

\section{Database Construction}
\subsection{Reference \& Distorted Image Collection}

To comprehensively characterize the data in Embodied scenarios, we collect 1,230 high-quality samples as reference images. (see Supplementary for data source) All data are pre-processed by Q-Align \cite{quality:q-align} to avoid pre-distortion before adding distortion. We focused on two aspects, Sim2Real and Perspective, to ensure coverage of both real and simulation, as well as first-person and third-person perspectives. In addition, we divided the subjects and backgrounds into five categories each, as shown in Figure \ref{fig:dataset}, to ensure involvement of each category and thus ensure the versatility of the database.

For the distorted images, according to the corruption caused by the current communication protocols, 30 distortion types are considered and classified into 7 categories: Blur, various types of unclear image; Luminance, global brightness changes; Chrominance, global color changes; Noise, random noise of different distributions; Compression, codec algorithm like JPEG; Spatial, local pixel-level changes; and others. For each distortion, we defined 5 intensity levels, ensuring the quality degradation perceived by the HVS is aligned at the same level. Thus, for each reference image, we randomly selected the intensity to add all the corruptions mentioned above, resulting in 36,900 distorted images. The reference/distorted image pairs will then be annotated by Embodied AI subjects.

\begin{table*}[t]
\centering
    \renewcommand\arraystretch{1.25}
    \renewcommand\tabcolsep{4pt}
    \belowrulesep=0pt\aboverulesep=0pt
    \caption{Comparison of Embodied-IQA with other perceptual quality databases. As a machine-oriented database, Embodied-IQA has not only more image samples and a larger annotation scale, but also comprehensive labels on three downstream steps. [Keys:  Cognition, Decision, Execution]}
    \label{tab:database}
    \resizebox{0.99\linewidth}{!}{
    \begin{tabular}{l|ccc|cc|cccccc}
    \toprule
    {\multirow{2}{*}{Database}}          & \multicolumn{3}{c|}{Image} & \multicolumn{2}{c|}{Corruption} & \multicolumn{6}{c}{Annotation}                          \\ \cdashline{2-12}
              & Reference  & Distorted  & Resolution  & Types        & Strength        & Num   & Dimension & Subjects & Cog. & Dec. & Exe.              \\ \midrule
    LIVE\cite{dataset:live}      & 29   & 779  & 768         & 5            & 5               & 25k   & 1   & Human (General) & {\faCheckSquare} &  &                   \\
    TID2013\cite{dataset:tid}       & 25   & 3k   & 512         & 24           & 5               & 514k  & 1   & Human (General) & {\faCheckSquare} &  &                   \\
    KADID-10K\cite{dataset:kadid10k} & 81   & 10k  & 1k          & 25           & 5               & 304k  & 1   & Human (General) & {\faCheckSquare} &  &                   \\
    CLIC2021\cite{dataset:clic}  & 585  & 3k   & 1k          & 10           & 3               & 484k  & 1   & Human (General) & {\faCheckSquare} &  &           \\
    NTIRE2022\cite{dataset:ntire2022} & 250  & 29k  & 288         & 40           & 5               & 1.13m & 1   & Human (General) & {\faCheckSquare} &  &       \\
    AGIQA-3K\cite{dataset:agiqa-3k}  & -    & 3k   & 1k          & -            & -               & 125k  & 1+1   & Human (Multimodal) & {\faCheckSquare} & {\faCheckSquare} &                  \\
    NTIRE2024\cite{dataset:aigiqa20k} & -    & 20k  & 1k          & -            & -               & 420k  & 1+1   & Human (Multimodal) & {\faCheckSquare} & {\faCheckSquare} &                  \\
    MPD \cite{dataset:mpd}     & 1k   & 30k  & 1k          & 30           & 5               & 2.25m & 5   & Machine (General) & {\faCheckSquare} &  &  \\ 
    EPD \cite{dataset:epd}     & 100   & 2.5k  & 256          & 25           & 5               & 30k & 2  & Machine (General) &  & {\faCheckSquare} &  \\ 
    \CLA{Embodied-IQA} {(ours)}       & \CLA{1.23k}   & \CLA{36.9k}  & \CLA{1k}          & \CLA{30}           & \CLA{5}               & \CLA{5.53m} & \CLA{3+3+1}   & \CLA{Machine (Robot)} & \CLA{\faCheckSquare} & \CLA{\faCheckSquare} & \CLA{\faCheckSquare} \\ \bottomrule
    \end{tabular}}
    \vspace{-4mm}
\end{table*}

\subsection{Perception: Task Definition}
\label{sec:task}

Perception refers to receiving information about the external environment, where Embodied AI obtains information through sensors, similar to human sensory organs. Considering that more than 82\% of human external input signals come from vision, we simplify this step of Embodied AI to the camera. First of all, we need to clarify the factors that Embodied AI focuses on in Perception. When viewing an image, HVS pays attention to factors such as brightness/chromaticity, while MVS/RVS relies on specific downstream tasks. Therefore, based on information such as objects, layout, and environment in the image, we manually annotate 5 tasks for each reference sample in natural language, as in previous MVS \cite{dataset:mpd} works. The difficulty of the tasks here increases in sequence and is limited to [\texttt{Cover}, \texttt{Insert}, \texttt{Move}, \texttt{Pick}, \texttt{Place}, \texttt{Pour}, \texttt{Press}, \texttt{Pull}, \texttt{Push}, \texttt{Twist}] to avoid being too difficult. All subsequent steps are based on the task corresponding to each image, and the image quality depends on the similarity of the inference results of the reference/distorted image pair.

\begin{figure}[t]
    \centering
    \includegraphics[width=\linewidth]{figs/dataset.pdf}
    \caption{Database construction of the Embodied-IQA, with 30k+ large-scale reference/distorted image pairs, meticulously annotated with 2m+ fine-grained \CLA{Cognition} score from 15 mainstream VLMs, 2m+ \CLA{Decision} score from 15 VLAs, and 1.5k real-world experiments as \CLA{Execution} score.}
    \label{fig:dataset}
    \vspace{-3mm}
\end{figure}

\subsection{Cognition: VLM Annotation}
\label{sec:vlm}

Cognition refers to the process of processing and understanding information after perception, including recognition, classification, memory, and reasoning. This function of Embodied AI is implemented through VLM, which corresponds to the human cerebrum. Specifically, we selected 15 commonly used VLMs for cognition, including: Mini-InternVL \cite{vlm:mini_internvl}, InternLM-Xcomposer2 \cite{vlm:internlm-xcomposer2}, InternLM-Xcomposer2.5 \cite{vlm:internlm-xcomposer2d5}, InternVL2 \cite{vlm:internvl2}, InternVL2.5 \cite{vlm:internvl2.5}, InternVL3 \cite{vlm:internvl3}, MPlugOwl3 \cite{vlm:mplug_owl}, Ovis1.5-Gemma \cite{vlm:ovis-gemma}, Ovis1.6-Llama \cite{vlm:ovis-llama}, Ovis2 \cite{vlm:ovis2}, Phi3-Vision \cite{vlm:phi3}, Phi3.5-Vision \cite{vlm:phi3.5}, Phi4-Multimodal \cite{vlm:phi4}, Qwen2-VL \cite{vlm:qwen2}, and Qwen2.5-VL \cite{vlm:qwen2.5}. To ensure usability in the Real-world, the parameter size we selected is all below 8B to ensure real-time inference.

Considering the output modality is textual, VLMs will be required to solve the pre-defined task in about 10 words, and the difference between reference/distorted output sentences will be measured. Specifically, the difference between the two output sentences includes three dimensions: accuracy, recall, and semantics, which are realized by the three classic indicators: BLEU \cite{eval:bleu}, ROUGE \cite{eval:cider}, and CIDEr \cite{eval:cider}. Since the maximum value of CIDEr is 10, we use a weighted average of 1:1:0.1 and report the sum of the five task results as the final Cognition score.

\subsection{Decision: VLA Annotation}
\label{sec:vla}

Decision refers to selecting the best course of action based on goals, rules, and experience, according to Cognition results. This function of Embodied AI is implemented through VLA, which corresponds to the human cerebellum. Specifically, we selected 15 commonly used VLA for Decision, including: CogACT \cite{vla:cogact}, Embodied-CoT \cite{vla:embodied-cot}, Octo \cite{vla:octo}, OpenVLA \cite{vla:openvla_v1}, OpenVLA-Libero \cite{vla:openvla_v3}, OpenVLA-Goal \cite{vla:openvla_v2}, OpenVLA-Libero-Object \cite{vla:openvla_v2}, OpenVLA-Libero-Spatial \cite{vla:openvla_v2}, Pi0-Aloha-Pen \cite{vla:pi0_v3}, Pi0-Aloha-Towel \cite{vla:pi0_v3}, Pi0-Aloha-Tupperware \cite{vla:pi0_v3}, Pi0-Base \cite{vla:pi0_v1}, Pi0-Droid \cite{vla:pi0_fast}, Pi0-Fast \cite{vla:pi0_fast}, and RT-X-1 \cite{vla:rt_x}. The parameter size of these VLA is also controlled at 8B, similar to VLM. 

Noted that since Embodied-IQA first introduced VLA into the IQA task, we define the quality of VLA as three dimensions. First, we parse the 7-DoF Pose\footnote{We will discard information beyond the above 7-DoF like depth, for alignment between VLAs.} output field. According to the mechanism of VLA, the first three represent position \footnote{For two-arm VLAs, we only select the arm with the larger movement range.} (translation of the operator along the three-dimensional coordinate system, in mm), the middle three represent rotation (rotation of the operator along the three-dimensional coordinate system, in rad), and the last one represents state (opening and closing of the operator, range [0-1]). The position score is based on the spatial distance of the coordinate points obtained from the reference/distorted image, the rotation score is based on the cosine similarity of directional vectors, and the last one is the absolute difference. The three dimensions are averaged after 0-1 normalization, and report the sum of the five task results as the final Decision score.

\subsection{Execution: Real-world Validation}

Execution refers to the process of transforming decisions into actual movements. This part of Embodied AI relies on specific actuators, corresponding to the human motor system. Based on kinematic statistics, the upper limbs dominate among all muscles and complete over 50\% of daily movements. Therefore, we use the robotic arm as the most representative actuator. Specifically, we execute tasks based on the inference results of VLA, with three scenarios: (1) Success: Directly assign score 100 to the sample; (2) Failure: Measure the Euclidean distance between the reference and distorted results based on the final pose of the actuator, and deduct points in centimeters; (3) Emergency stop: If the actuator hits the table or wall, directly assign score 0. Considering the uncontrollable factors in real-machine experiments, we only execute the task with the lowest difficulty level among the 5 tasks to verify whether the results of VLM and VLA align with the Real-world.

\begin{figure*}[t]
    \centering
    \subfigure{\includegraphics[width=0.48\textwidth]{figs/legend-vlm.pdf}}\hskip.2em
    \subfigure{\includegraphics[width=0.48\textwidth]{figs/legend-vla.pdf}}\hskip.2em
    \vspace{-8pt}
    \caption{Benchmarking \CLB{VLMs}\&\CLA{VLAs} in 3 different score dimensions and 5 distortion levels. Their performance varies in 3 dimensions and decreases with the distortion. (Zoom in for detail)}
    \label{fig:radar}
    \vspace{-3mm}
\end{figure*}

\begin{figure*}[t]
    \centering
    \begin{minipage}[]{0.16\linewidth}
    \centering
    \centerline{\includegraphics[width = \textwidth]{figs/CorrVLM_bleu_0.3054.pdf}}
    \centerline{Precision (0.31)}\medskip
    \end{minipage}
    \begin{minipage}[]{0.16\linewidth}
    \centering
    \centerline{\includegraphics[width = \textwidth]{figs/CorrVLM_cider_0.3092.pdf}}
    \centerline{Recall (0.31)}\medskip
    \end{minipage}
    \begin{minipage}[]{0.16\linewidth}
    \centering
    \centerline{\includegraphics[width = \textwidth]{figs/CorrVLM_rouge_0.3042.pdf}}
    \centerline{Semantic (0.30)}\medskip
    \end{minipage}
    \begin{minipage}[]{0.16\linewidth}
    \centering
    \centerline{\includegraphics[width = \textwidth]{figs/CorrVLA_score1_0.2487.pdf}}
    \centerline{Position (0.25)}\medskip
    \end{minipage}
    \begin{minipage}[]{0.16\linewidth}
    \centering
    \centerline{\includegraphics[width = \textwidth]{figs/CorrVLA_score2_0.2319.pdf}}
    \centerline{Rotation (0.23)}\medskip
    \end{minipage}
    \begin{minipage}[]{0.16\linewidth}
    \centering
    \centerline{\includegraphics[width = \textwidth]{figs/CorrVLA_score3_0.2805.pdf}}
    \centerline{State (0.28)}\medskip
    \end{minipage}
    \vspace{-3mm}

    \caption{Correlation matrix of \CLB{VLMs}\&\CLA{VLAs} subjects, the a-o order follows Section \ref{sec:vlm},\ref{sec:vla}. Darker colors denote a higher SRCC, with the averaged SRCC attached to the bottom of the matrix.}
    \label{fig:corr}
    \vspace{-4mm}
\end{figure*}

\section{Database Analysis}

This section analyzes Embodied-IQA database from four dimensions: On the model level, we (1) Benchmark the VLMs and VLAs when processing the distorted images; (2) Explore the internal correlation between VLMs and VLAs; On the instance level, we (3) Compare the score distributions under different categories and distortions; (4) Analyze the distortion sensitivity of VLMs/VLAs.

\textbf{[Benchmark]} We select 6 representative VLM and VLA in Figure \ref{fig:radar}. As the distortion level increases, the total scores of both VLM and VLA gradually decrease. However, the differences among the three scoring dimensions of VLM are much greater than the level of distortion. After distortion, the Semantic score of the image decreases relatively little, followed by Recall, and then Precision. In VLM, the reference/distorted output of MPlugOwl3 is the most consistent, while advanced models like Qwen2.5-VL are less robust. Therefore, distortion usually affects VLM at the character level rather than the semantic level, and it is more likely to output redundant text than to lose information. Meanwhile, there are also significant differences among the three scoring dimensions of VLA. In VLA, Octo shows strong robustness to distortion in Position and Rotation, while models like CogACT and OpenVLA are more faithful in State. Among them, State changes little after distortion, Position changes more, and Rotation is the most easily affected by distortion. This indicates that distortion usually does not affect the end operator but has a significant impact on the robot arm. Therefore, VLA needs to be carefully selected to ensure high-quality output of the first 6-DoF.

\begin{figure}[t]
    \centering
    \includegraphics[width=\linewidth]{figs/main-subjective.pdf}
    \caption{Decision score visualized in 30 distortion subsets. Different color denotes distortion \GroupA{Level 1}-\GroupB{Level 2}-\GroupC{Level 3}-\GroupD{Level 4}-\GroupE{Level 5}. Different distortions affecting VLAs vary significantly.}
    \label{fig:subjective}
    \vspace{-2mm}


\end{figure}

\textbf{[Correlation]} We visualize the correlation between subject models in Figure \ref{fig:corr} according to the order of models in Section \ref{sec:vlm},\ref{sec:vla} through Spearman Rank-order Correlation Coefficient (SRCC). According to past IQA \cite{dataset:mpd} work, the correlation of HVS is usually above 0.6, while that of MVS is often less than 0.5. In specific machine tasks, detection has better correlation, while Visual Question Answering (VQA) may even be less than 0.4. For Cognition, which uses VLM to solve embodied tasks, is a degraded form of VQA, a correlation slightly higher than 0.3 is understandable. Moreover, we find that the correlation of VLA is even lower than that of VLM, at around 0.25. Therefore, when evaluating VLM and VLA, using only one model as the subject is far from sufficient, especially for VLA. It is necessary to collect their general preference. This also reflects the necessity of constructing the Embodied-IQA database and the separation of Cognition and Decision in the RVS.

\textbf{[Distribution]} Since Decision is more downstream than Cognition and has never been deeply investigated in IQA, we show the distribution of Decision and put Cognition in the supplementary. Figure \ref{fig:subjective} shows the Decision score distribution under 30 types of distortions and 5 intensity levels. Results show that RVS and the traditional HVS have significant differences. Taking brightness as an example, Embodied AI is highly sensitive to `Maximum brighten/darken', and the quality significantly decreases with the distortion level; however, it is rarely affected by `Mean brighten/darken', and there is no significant distribution change from level 1 to 5. These findings emphasize the differences between RVS and HVS. Figure \ref{fig:violin} lists the relation of the three Decision dimensions and the score distributions corresponding to different image categories. Overall, Position, Rotation, and State show independent distributions. In the Sim2Real distribution, the real scores are relatively high, indicating that VLA is better at handling real-world data; the first-person results are far worse than the third-person results, indicating that in the training data of VLA, the sampling tools and actuators are rarely integrated, which needs to be improved in the future; there is also a certain fluctuation between the five Main Objects and Background. These two findings jointly support the rationality of the division of source image data and annotation dimensions in Embodied-IQA.

\textbf{[Sensitivity]} As mentioned in the Distribution section, the difference between MVS and HVS causes VLM to be highly sensitive to some distortion categories, but robust to others; the difference between RVS and HVS also causes VLA to have a similar phenomenon. Therefore, we combined the Just-Noticable-Difference (JND) theory to analyze the similarities and differences in the distortion sensitivity of VLM and VLA, as shown in Figure \ref{fig:jnd}. The sensitivity is divided into three levels, Mild, Medium, and Severe, each accounting for one-third of all image samples, according to the Cognition and Decision scores. Results for VLM and VLA are significantly different from human perception. For example, Dis 16 (Gaussian denoise) will cause serious distortion at level 1, and Dis25 (Block interpolation) has no significant effect on VLM/VLA even at level 5. Meanwhile, although VLM and VLA have certain commonalities in sensitivity, there are also distortions such as Dis02 (Lens blur) that mainly affect VLM, or Dis15 (Multiplicative noise) that mainly affect VLA. This VLM\&VLA-based partition further explains the gap between Cognition and Decision, and can serve as an important reference in the IQA for the Embodied AI topic in the future.

\begin{figure*}[t]
    \centering
    \begin{minipage}[]{0.21\linewidth}
    \centering
    \centerline{\includegraphics[width = \textwidth]{figs/DimVLA.pdf}}
    % \centerline{VQA}\medskip
    \end{minipage}
    \begin{minipage}[]{0.75\linewidth}
    \centering
    \centerline{\includegraphics[width = \textwidth]{figs/vla-subtypes.pdf}}
    % \centerline{CAP}\medskip
    \end{minipage}
    \vspace{-2mm}

    \caption{Correlation between the general Decision score and the 3 dimensions from VLAs, and the score distribution in Sim2Real, First/Third perspectives, Main object, and Background sub-categories.}
    \label{fig:violin}
    \vspace{-1mm}
\end{figure*}

\begin{figure*}[t]
    \centering
    \begin{minipage}[]{0.86\linewidth}
    \centering
    \centerline{\includegraphics[width = \textwidth]{figs/jnd-main.pdf}}
    % \centerline{VQA}\medskip
    \end{minipage}
    \begin{minipage}[]{0.12\linewidth}
    \centering
    \centerline{\includegraphics[width = \textwidth]{figs/jnd-label.pdf}}
    % \centerline{CAP}\medskip
    \end{minipage}
    \vspace{-2mm}

    \caption{Just-Noticable-Difference (JND) of \CLB{VLMs} and \CLA{VLAs}. Distortion order follows Figure \ref{fig:subjective}.}
    \label{fig:jnd}
    \vspace{-1mm}
\end{figure*}

\section{Experiment}

\subsection{Experiment Setups}
\label{sec:train}

We randomly partitioned the Embodied-IQA database into the train/val set, with 29,520 and 7,380 reference/distorted image pairs according to an 8:2 ratio. 
Since VLA is more downstream than VLM in the Mertonian system of Embodied AI, and previous machine-oriented IQA works are all based on VLM, we use the Decision scores from VLA in the main experiment and put the Cognition scores from VLM in the supplementary.
15 objective IQA metrics are implemented to predict the Decision score, including:
(1) 5 baseline metrics: PSNR, SSIM \cite{quality:ssim}, Brisque \cite{quality:brisque}, Q-Align \cite{quality:q-align}, and Q-Align+ \cite{quality:q-align+}. Q-Align and Q-Align+ are loaded in quality/aesthetic weights. As the most commonly used IQA metrics, they are performed in a Zero-shot setting; 
(2) 5 Full-Reference (FR) metrics: AHIQ \cite{quality:ahiq}, CKDN \cite{quality:ckdn}, DISTS \cite{quality:dists}, LPIPS \cite{quality:lpips}, and TOPIQ-FR \cite{quality:topiq}. Those learning-based metrics are fine-tuned on our training set, which takes both reference/distorted images as inputs;
(3) 5 No-Reference (NR) metrics: CLIPIQA \cite{quality:CLIPIQA}, CNNIQA \cite{quality:CNNIQA}, DBCNN \cite{quality:DBCNN}, QualiClip \cite{quality:qualiclip}, and TOPIQ-NR \cite{quality:topiq}. They are also fine-tuned on a training set, which takes only distorted images as inputs.

To benchmark the performance of quality metrics, three global indicators were employed: SRCC, Kendall Rank-order Correlation Coefficient (KRCC), and Pearson Linear Correlation Coefficient (PLCC), to evaluate the consistency between the objective quality score and the subjective MOS. Among these, SRCC and KRCC represent the prediction monotonicity, while PLCC measures the accuracy.
We train the FR/NR metrics on Embodied-IQA with the learning rate as $10^{-5}$ for 50 epochs, under the default settings in pyiqa toolbox, and evaluate the performance on (1) 3 scoring dimension; (2) 3 JND-based distortion sensitivity in Figure \ref{fig:jnd}; (3) First/Third person perspective; (4) Sim2Real; (5) 5 Distortion level. 
The partitioning, training, and testing pipeline is repeated 10 times, and the mean value is reported as the experimental result.
The perception module is based on the Intel RealSense D455 array, supporting both First-person (wrist) and Third-person (top, side) as input.
Cognition and Decision annotations are collected on two servers with 16$\times$NVIDIA A800 SXM4 80GB GPUs, and then conduct IQA training/validation on one GPU above.
Execution is achieved through the UR5 robotic arm and Robotiq 2F-140 gripper, with a working radius of 85cm.


\subsection{Result and Discussion}

Table \ref{tab:exp-content} presents the performance of advanced quality metrics on the Embodied-IQA database.
For the three dimensions of the 7-DoF VLA output, Position is the easiest to predict, followed by State, with Rotation being the most difficult. The SRCC of FR IQA methods with subjective labels is less than 0.65, while NR is even less than 0.6. Note that these methods have a correlation close to 0.9 with HVS in traditional human-oriented IQA tasks, which is sufficiently excellent, but they cannot adapt to the database we proposed, indicating that IQA for Embodied AI needs further research. For distortion sensitivity, the Decision scores are high under mild distortions, with small internal gaps and difficulty in prediction; whereas when the distortion becomes severe, the Distortion scores fluctuate more, resulting in a higher SRCC. For Perspective and Sim2Real, most IQA methods perform better on third-person, real images. Therefore, in Embodied scenarios, more content captured by the robotic arm itself or from simulation software should be used. For the five absolute distortion levels, the performance of IQA methods does not change much. This further proves that dividing distortion levels based on HVS is insufficient, and the distortion levels of Embodied AI should be divided by the JND of RVS itself, as we have done in the Embodied-IQA database. Comparing various IQA methods longitudinally, FR is superior to NR in most cases, with TOPIQ maintaining the leading performance in most cases, with an SRCC of about 0.75, which still needs improvement compared to human-oriented IQA. It is worth mentioning that the main parameters of some methods have been frozen based on HVS, such as LPIPS, DISTS, and CLIPIQA. Thus, although after training, they are even worse than the zero-shot baseline. This further reflects the gap between HVS and RVS, implying the significance of proposing the Embodied IQA task.

\begin{table*}[t]
    \centering
    \renewcommand\arraystretch{1.25}
    \renewcommand\tabcolsep{2.7pt}
    \belowrulesep=0pt\aboverulesep=0pt
    \caption{Using 15 advanced IQA metrics to evaluate the \CLA{Decision} score from VLAs, including zero-shot, FR, and NR metrics. [Keys: \CLA{Best}/\CLB{Second best} in group; \textbf{Baseline}; \colorbox{gray!20}{Lower} than baseline.]}
    \label{tab:exp-content}

    \resizebox{\linewidth}{!}{
\begin{tabular}{l|l|ccc|ccc|ccc|ccc|ccc}
\toprule
\multicolumn{2}{c|}{Dimension}            & \multicolumn{3}{c|}{Position} & \multicolumn{3}{c|}{Rotation}  & \multicolumn{3}{c|}{State}  & \multicolumn{3}{c|}{First Perspective}& \multicolumn{3}{c}{Third Perspective}   \\ \cline{1-17}
Group                 & \multicolumn{1}{c|}{Metric}        & SRCC$\uparrow$   & KRCC$\uparrow$   & PLCC$\uparrow$   & SRCC$\uparrow$   & KRCC$\uparrow$   & PLCC$\uparrow$   & SRCC$\uparrow$   & KRCC$\uparrow$   & PLCC$\uparrow$   & SRCC$\uparrow$   & KRCC$\uparrow$   & PLCC$\uparrow$   & SRCC$\uparrow$   & KRCC$\uparrow$   & PLCC$\uparrow$   \\ \midrule
\multirow{5}{*}{Zero} & PSNR          &0.2762&0.1872&0.3094&0.2594&0.1756&0.3035&0.2284&0.1522&0.2271&0.4059&0.2763&0.4373&0.3949&0.2693&0.4376


 \\
                      & SSIM \cite{quality:ssim}         & 0.4862&0.3345&0.4438&0.4246&0.2891&0.3912&0.3607&0.2468&0.3216&0.5834&0.4101&0.5256&0.5478&0.3815&0.5132


 \\
                      & Brisque \cite{quality:brisque}      & 0.3073&0.2051&0.2707&0.2752&0.1826&0.2519&0.3335&0.2255&0.2986&0.3302&0.2210&0.2634&0.4020&0.2688&0.3836


 \\
                      & Q-Align \cite{quality:q-align}  &\textbf{0.5325}&\textbf{0.3641}&\textbf{0.4869}&\textbf{0.5387}&\textbf{0.3758}&\textbf{0.5329}&\textbf{0.3791}&\textbf{0.2552}&\textbf{0.3346}&\textbf{0.6658}&\textbf{0.4715}&\textbf{0.5992}&\textbf{0.5854}&\textbf{0.4030}&\textbf{0.5578}
\\
                      & Q-Align+ \cite{quality:q-align+} & 0.3275&0.2157&0.2410&0.3596&0.2465&0.2818&0.1492&0.0972&0.1272&0.4663&0.3167&0.4283&0.3104&0.2021&0.2649

 \\ \cdashline{1-17}
\multirow{5}{*}{FR}   & AHIQ \cite{quality:ahiq}         & \CLB{0.7481}&\CLB{0.5496}&\CLB{0.7467}&\CLA{0.6454}&\CLA{0.4655}&\CLB{0.6435}&\CLB{0.6465}&\CLB{0.4609}&\CLB{0.6590}&\CLB{0.8011}&\CLB{0.6014}&\CLB{0.8025}&\CLB{0.7989}&\CLB{0.6007}&\CLB{0.7959}
\\
                      & CKDN \cite{quality:ckdn}         & 0.6748&0.4807&0.6771&0.6061&0.4278&0.6001&0.6324&0.4515&0.6410&0.7716&0.5720&0.7610&0.7641&0.5624&0.7596
\\
                      & DISTS \cite{quality:dists}        & 0.5797&0.4010&0.5846&\colorbox{gray!20}{0.5366}&\colorbox{gray!20}{0.3746}&0.5390&0.4653&0.3180&0.4611&\colorbox{gray!20}{0.6458}&\colorbox{gray!20}{0.4624}&0.6249&0.6545&0.4654&0.6642
 \\
                      & LPIPS \cite{quality:lpips}        & \colorbox{gray!20}{0.3922}&\colorbox{gray!20}{0.2642}&\colorbox{gray!20}{0.3511}&\colorbox{gray!20}{0.2972}&\colorbox{gray!20}{0.1994}&\colorbox{gray!20}{0.2378}&0.4210&0.2890&0.3852&\colorbox{gray!20}{0.4697}&\colorbox{gray!20}{0.3205}&\colorbox{gray!20}{0.4168}&\colorbox{gray!20}{0.4821}&\colorbox{gray!20}{0.3313}&\colorbox{gray!20}{0.4805}
 \\
                      & TOPIQ-FR \cite{quality:topiq}     & \CLA{0.7748}&\CLA{0.5794}&\CLA{0.7827}&\CLB{0.6428}&\CLB{0.4607}&\CLA{0.6480}&\CLA{0.6684}&\CLA{0.4826}&\CLA{0.6727}&\CLA{0.8307}&\CLA{0.6371}&\CLA{0.8297}&\CLA{0.8322}&\CLA{0.6404}&\CLA{0.8298}
\\ \cdashline{1-17}
\multirow{6}{*}{NR}   & CLIPIQA \cite{quality:CLIPIQA}      & \colorbox{gray!20}{0.1784}&\colorbox{gray!20}{0.1172}&\colorbox{gray!20}{0.1246}&\colorbox{gray!20}{0.0708}&\colorbox{gray!20}{0.0193}&\colorbox{gray!20}{0.0468}&\colorbox{gray!20}{0.1348}&\colorbox{gray!20}{0.0770}&\colorbox{gray!20}{0.0622}&\colorbox{gray!20}{0.0048}&\colorbox{gray!20}{0.0043}&\colorbox{gray!20}{0.0821}&\colorbox{gray!20}{0.2155}&\colorbox{gray!20}{0.1287}&\colorbox{gray!20}{0.1415}

 \\
                      & CNNIQA \cite{quality:CNNIQA}       & \colorbox{gray!20}{0.5189}&0.3642&0.5318&\colorbox{gray!20}{0.4618}&\colorbox{gray!20}{0.3221}&\colorbox{gray!20}{0.4587}&0.4601&0.3203&0.4667&\colorbox{gray!20}{0.5441}&\colorbox{gray!20}{0.3770}&\colorbox{gray!20}{0.5618}&0.6407&0.4579&0.6468

\\
                      & DBCNN \cite{quality:DBCNN}        & 0.6045&0.4303&0.6094&\colorbox{gray!20}{0.5325}&\colorbox{gray!20}{0.3687}&0.5341&0.5419&0.3761&0.5441&\colorbox{gray!20}{\CLB{0.6399}}&\colorbox{gray!20}{\CLB{0.4593}}&\CLB{0.6408}&0.6565&0.4653&0.6609

 \\
                      & QualiClip \cite{quality:qualiclip}   & \CLB{0.6463}&\CLB{0.4619}&\CLB{0.6643}&\CLB{0.5387}&\CLB{0.3768}&\CLB{0.5384}&\CLB{0.5589}&\CLB{0.3912}&\CLB{0.5518}&\colorbox{gray!20}{0.5428}&\colorbox{gray!20}{0.3870}&\colorbox{gray!20}{0.5490}&\CLB{0.7208}&\CLB{0.5240}&\CLB{0.7175}

\\
                      & TOPIQ-NR \cite{quality:topiq}     & \CLA{0.7496}&\CLA{0.5549}&\CLA{0.7606}&\CLA{0.5981}&\CLA{0.4253}&\CLA{0.6020}&\CLA{0.7036}&\CLA{0.5100}&\CLA{0.6960}&\CLA{0.7791}&\CLA{0.5804}&\CLA{0.7810}&\CLA{0.8269}&\CLA{0.6341}&\CLA{0.8211} \\ \bottomrule
\end{tabular}}
\resizebox{\linewidth}{!}{
\begin{tabular}{l|l|ccc|ccc|ccc|ccc|ccc}
\toprule
\multicolumn{2}{c|}{Dimension}          & \multicolumn{3}{c|}{Mild Distortion}  & \multicolumn{3}{c|}{Medium Distortion} & \multicolumn{3}{c|}{Severe Distortion}  & \multicolumn{3}{c|}{Real-world} & \multicolumn{3}{c}{Simulation}   \\ \cline{1-17}
Group                 & \multicolumn{1}{c|}{Metric}        & SRCC$\uparrow$   & KRCC$\uparrow$   & PLCC$\uparrow$   & SRCC$\uparrow$   & KRCC$\uparrow$   & PLCC$\uparrow$   & SRCC$\uparrow$   & KRCC$\uparrow$   & PLCC$\uparrow$   & SRCC$\uparrow$   & KRCC$\uparrow$   & PLCC$\uparrow$   & SRCC$\uparrow$   & KRCC$\uparrow$   & PLCC$\uparrow$   \\ \midrule
\multirow{5}{*}{Zero} & PSNR          & 0.3518&0.2396&\textbf{0.4935}&0.2292&0.1555&0.2175&0.1190&0.0806&0.1443&0.3794&0.2575&0.3767&0.2617&0.1778&0.3905



 \\
                      & SSIM \cite{quality:ssim}         & \textbf{0.3778}&\textbf{0.2620}&0.2581&0.1993&0.1333&0.1617&0.2387&0.1604&0.2754&0.5405&0.3754&0.4848&0.5336&0.3724&0.4990



 \\
                      & Brisque \cite{quality:brisque}      & 0.2088&0.1334&0.1502&0.1503&0.0967&0.1168&0.1525&0.1017&0.1205&0.4143&0.2798&0.3907&0.0636&0.0421&0.0423



 \\
                      & Q-Align \cite{quality:q-align} & 0.3541&0.2403&0.3097&\textbf{0.4090}&\textbf{0.2873}&\textbf{0.3899}&\textbf{0.3258}&\textbf{0.2206}&\textbf{0.3723}&\textbf{0.6179}&\textbf{0.4302}&\textbf{0.5639}&\textbf{0.6565}&\textbf{0.4690}&\textbf{0.6333}
 \\
                      & Q-Align+ \cite{quality:q-align+} & 0.1699&0.1132&0.1375&0.2799&0.1908&0.2287&0.0618&0.0425&0.1506&0.4167&0.2753&0.3820&0.4846&0.3285&0.4711

 \\ \cdashline{1-17}
\multirow{5}{*}{FR}   & AHIQ \cite{quality:ahiq}         & \CLB{0.6683}&\CLB{0.4831}&\CLB{0.6932}&\CLB{0.6653}&0.4821&\CLB{0.7130}&\CLB{0.7218}&\CLB{0.5308}&\CLB{0.7278}&\CLA{0.8138}&\CLA{0.6223}&\CLA{0.8270}&0.7515&0.5583&0.7520

 \\
                      & CKDN \cite{quality:ckdn}         & 0.5733&0.4060&0.5743&0.6610&\CLB{0.4841}&0.6979&0.6909&0.5024&0.6867&0.7227&0.5303&0.7311&\CLB{0.7676}&\CLB{0.5769}&\CLB{0.7769}

\\
                      & DISTS \cite{quality:dists}        & 0.4564&0.3191&\colorbox{gray!20}{0.4507}&\colorbox{gray!20}{0.2343}&\colorbox{gray!20}{0.1604}&\colorbox{gray!20}{0.2735}&\colorbox{gray!20}{0.3177}&\colorbox{gray!20}{0.2155}&0.3731&0.6443&0.4551&0.6394&\colorbox{gray!20}{0.6408}&\colorbox{gray!20}{0.4589}&0.6560

\\
                      & LPIPS \cite{quality:lpips}        & \colorbox{gray!20}{0.3004}&\colorbox{gray!20}{0.1986}&\colorbox{gray!20}{0.2102}&0.4208&0.2884&0.4581&0.4672&0.3220&0.4719&\colorbox{gray!20}{0.3605}&\colorbox{gray!20}{0.2429}&\colorbox{gray!20}{0.3999}&\colorbox{gray!20}{0.4665}&\colorbox{gray!20}{0.3181}&\colorbox{gray!20}{0.4442}

 \\
                      & TOPIQ-FR \cite{quality:topiq}     & \CLA{0.7128}&\CLA{0.5257}&\CLA{0.7626}&\CLA{0.7238}&\CLA{0.5371}&\CLA{0.7581}&\CLA{0.7355}&\CLA{0.5434}&\CLA{0.7434}&\CLB{0.8104}&\CLB{0.6210}&\CLB{0.8250}&\CLA{0.7815}&\CLA{0.5910}&\CLA{0.8053}

 \\ \cdashline{1-17}
\multirow{6}{*}{NR}   & CLIPIQA \cite{quality:CLIPIQA}      & \colorbox{gray!20}{0.0848}&\colorbox{gray!20}{0.0563}&\colorbox{gray!20}{0.0568}&\colorbox{gray!20}{0.0413}&\colorbox{gray!20}{0.0300}&\colorbox{gray!20}{0.0779}&\colorbox{gray!20}{0.2198}&\colorbox{gray!20}{0.1509}&\colorbox{gray!20}{0.2150}&\colorbox{gray!20}{0.1576}&\colorbox{gray!20}{0.1040}&\colorbox{gray!20}{0.1293}&\colorbox{gray!20}{0.1298}&\colorbox{gray!20}{0.0863}&\colorbox{gray!20}{0.1093}

 \\
                      & CNNIQA \cite{quality:CNNIQA}       & \colorbox{gray!20}{0.3766}&\colorbox{gray!20}{0.2607}&0.3847&\colorbox{gray!20}{0.3367}&\colorbox{gray!20}{0.2314}&0.3907&0.3921&0.2696&0.3957&\colorbox{gray!20}{0.5651}&\colorbox{gray!20}{0.3967}&0.5875&\colorbox{gray!20}{0.5835}&\colorbox{gray!20}{0.4117}&\colorbox{gray!20}{0.5901}

\\
                      & DBCNN \cite{quality:DBCNN}        & 0.4488&0.3056&0.4309&0.4283&0.2994&\CLB{0.4699}&\CLB{0.4622}&\CLB{0.3195}&\CLB{0.4459}&\CLB{0.6728}&0.4833&0.6806&\colorbox{gray!20}{\CLB{0.6468}}&\colorbox{gray!20}{\CLB{0.4613}}&\CLB{0.6345}

 \\
                      & QualiClip \cite{quality:qualiclip}   &\CLB{0.4941}&\CLB{0.3425}&\CLB{0.4794}&\CLB{0.4390}&\CLB{0.3094}&0.4607&0.3752&0.2633&0.3917&0.6712&\CLB{0.4882}&\CLB{0.6952}&\colorbox{gray!20}{0.6308}&\colorbox{gray!20}{0.4468}&\colorbox{gray!20}{0.6292}

 \\
                      & TOPIQ-NR \cite{quality:topiq}     & \CLA{0.7035}&\CLA{0.5164}&\CLA{0.7263}&\CLA{0.7174}&\CLA{0.5312}&\CLA{0.7374}&\CLA{0.7227}&\CLA{0.5312}&\CLA{0.7310}&\CLA{0.7995}&\CLA{0.6072}&\CLA{0.8148}&\CLA{0.7697}&\CLA{0.5771}&\CLA{0.7777}

 \\ \bottomrule
\end{tabular}}
\resizebox{\linewidth}{!}{
\begin{tabular}{l|l|ccc|ccc|ccc|ccc|ccc}
\toprule
\multicolumn{2}{c|}{Dimension}           & \multicolumn{3}{c|}{Dis-level-1}  & \multicolumn{3}{c|}{Dis-level-2} & \multicolumn{3}{c|}{Dis-level-3}  & \multicolumn{3}{c|}{Dis-level-4}  & \multicolumn{3}{c}{Dis-level-5}  \\ \cline{1-17}
Group                 & \multicolumn{1}{c|}{Metric}        & SRCC$\uparrow$   & KRCC$\uparrow$   & PLCC$\uparrow$   & SRCC$\uparrow$   & KRCC$\uparrow$   & PLCC$\uparrow$   & SRCC$\uparrow$   & KRCC$\uparrow$   & PLCC$\uparrow$   & SRCC$\uparrow$   & KRCC$\uparrow$   & PLCC$\uparrow$   & SRCC$\uparrow$   & KRCC$\uparrow$   & PLCC$\uparrow$   \\ \midrule
\multirow{5}{*}{Zero} & PSNR          & 0.2932&0.1987&0.3856&0.2322&0.1532&0.2837&0.2215&0.1484&0.2749&0.2379&0.1591&0.2642&0.3575&0.2444&0.3571


 \\
                      & SSIM \cite{quality:ssim}         & 0.4397&0.3035&0.4038&0.4638&0.3198&0.4208&0.4865&0.3329&0.4446&0.4927&0.3383&0.4745&0.5558&\textbf{0.3832}&\textbf{0.5327}


 \\
                      & Brisque \cite{quality:brisque}      & 0.2650&0.1786&0.2472&0.2446&0.1618&0.2443&0.3516&0.2346&0.3033&0.2928&0.1933&0.2604&0.3632&0.2434&0.3269


 \\
                      & Q-Align \cite{quality:q-align}  & \textbf{0.5049}&\textbf{0.3488}&\textbf{0.4924}&\textbf{0.5534}&\textbf{0.3865}&\textbf{0.4951}&\textbf{0.5609}&\textbf{0.3876}&\textbf{0.5132}&\textbf{0.5753}&\textbf{0.4015}&\textbf{0.5124}&\textbf{0.5567}&0.3818&0.5275
 \\
                      & Q-Align+ \cite{quality:q-align+} & 0.2909&0.1923&0.2210&0.4074&0.2763&0.3245&0.3828&0.2537&0.3241&0.3862&0.2585&0.3145&0.3274&0.2200&0.2694

 \\ \cdashline{1-17}
\multirow{5}{*}{FR}   & AHIQ \cite{quality:ahiq}         & \CLB{0.7453}&\CLB{0.5531}&\CLB{0.7722}&\CLB{0.7610}&\CLB{0.5671}&\CLB{0.7741}&\CLB{0.7389}&\CLB{0.5397}&\CLB{0.7546}&\CLB{0.7486}&\CLB{0.5518}&\CLB{0.7692}&\CLB{0.7655}&\CLB{0.5698}&\CLB{0.7820}

 \\
                      & CKDN \cite{quality:ckdn}         & 0.6253&0.4513&0.6806&0.6612&0.4785&0.6888&0.7036&0.5106&0.7124&0.7030&0.5131&0.7035&0.7194&0.5271&0.7290

 \\
                      & DISTS \cite{quality:dists}        & 0.5659&0.3962&0.5710&0.5774&0.4079&0.5753&0.5611&0.3903&0.5636&\colorbox{gray!20}{0.5431}&\colorbox{gray!20}{0.3766}&0.5727&0.5834&0.4015&0.5967

 \\
                      & LPIPS \cite{quality:lpips}        & \colorbox{gray!20}{0.2736}&\colorbox{gray!20}{0.1810}&\colorbox{gray!20}{0.2813}&\colorbox{gray!20}{0.3723}&\colorbox{gray!20}{0.2550}&\colorbox{gray!20}{0.4003}&\colorbox{gray!20}{0.4354}&\colorbox{gray!20}{0.2927}&\colorbox{gray!20}{0.4208}&\colorbox{gray!20}{0.4619}&\colorbox{gray!20}{0.3156}&\colorbox{gray!20}{0.4567}&\colorbox{gray!20}{0.4824}&\colorbox{gray!20}{0.3322}&\colorbox{gray!20}{0.4526}

 \\
                      & TOPIQ-FR \cite{quality:topiq}     & \CLA{0.7861}&\CLA{0.5939}&\CLA{0.8034}&\CLA{0.7765}&\CLA{0.5828}&\CLA{0.8012}&\CLA{0.7741}&\CLA{0.5809}&\CLA{0.7871}&\CLA{0.7629}&\CLA{0.5682}&\CLA{0.7773}&\CLA{0.7763}&\CLA{0.5766}&\CLA{0.7873}

 \\ \cdashline{1-17}
\multirow{6}{*}{NR}   & CLIPIQA \cite{quality:CLIPIQA}       & \colorbox{gray!20}{0.1541}&\colorbox{gray!20}{0.1054}&\colorbox{gray!20}{0.0822}&\colorbox{gray!20}{0.1204}&\colorbox{gray!20}{0.0772}&\colorbox{gray!20}{0.0635}&\colorbox{gray!20}{0.1291}&\colorbox{gray!20}{0.0831}&\colorbox{gray!20}{0.0726}&\colorbox{gray!20}{0.1564}&\colorbox{gray!20}{0.1041}&\colorbox{gray!20}{0.1129}&\colorbox{gray!20}{0.1130}&\colorbox{gray!20}{0.0732}&\colorbox{gray!20}{0.0846}

 \\
                      & CNNIQA \cite{quality:CNNIQA}       & \colorbox{gray!20}{0.4871}&\colorbox{gray!20}{0.3362}&0.5207&\colorbox{gray!20}{0.5215}&\colorbox{gray!20}{0.3669}&0.5522&\colorbox{gray!20}{0.5441}&\colorbox{gray!20}{0.3835}&0.5710&\colorbox{gray!20}{0.5433}&\colorbox{gray!20}{0.3787}&0.5570&0.6083&0.4274&0.6078

 \\
                      & DBCNN \cite{quality:DBCNN}        & \CLB{0.5982}&\CLB{0.4225}&\CLB{0.6213}&\CLB{0.6194}&\CLB{0.4431}&0.6210&\CLB{0.6235}&\CLB{0.4386}&\CLB{0.6146}&0.6189&0.4394&0.6163&\CLB{0.6141}&\CLB{0.4335}&0.6282

\\
                      & QualiClip \cite{quality:qualiclip} &0.5805&0.4137&0.5787&0.6029&0.4342&\CLB{0.6389}&0.6044&0.4372&0.6101&\CLB{0.6333}&\CLB{0.4509}&\CLB{0.6469}&0.6030&0.4265&\CLB{0.6399}

 \\
                      & TOPIQ-NR \cite{quality:topiq}     & \CLA{0.7547}&\CLA{0.5626}&\CLA{0.7781}&\CLA{0.7868}&\CLA{0.5937}&\CLA{0.7989}&\CLA{0.7566}&\CLA{0.5601}&\CLA{0.7695}&\CLA{0.7480}&\CLA{0.5532}&\CLA{0.7658}&\CLA{0.7529}&\CLA{0.5558}&\CLA{0.7670}

 \\ \bottomrule
\end{tabular}}
\vspace{-4mm}
\end{table*}

\subsection{Cross Database Validation}

To further analyze the differences between evaluating human and Embodied AI preference, we conducted cross-validation using Embodied-IQA VLA Decision score for training, VLM Cognition score, and two of the most commonly human-oriented databases, LIVE2008 \cite{dataset:live} and TID2013 \cite{dataset:tid} for validation, employing the same parameter settings as the Section \ref{sec:train}.
According to the results in Table \ref{fig:subject-bar} (a), we find that IQA methods fine-tuned on Embodied AI data lose certain human-oriented evaluation capabilities, where the SRCC is even lower than 0.4 in the LIVE database. 
Fortunately, the IQA model trained with VLA annotations can also predict VLM scores, and the SRCC of AHIQ can reach 0.7, revealing the internal connection between Cognition and Decision.

\subsection{Real-world Experiment}

Since Embodied AI is ultimately applied in the Real-world, we compare Execution with Cognition/Decision to link External and Internal Reality, thereby proving the reliability of the 5m+ annotations in the Embodied-IQA database. Specifically, we selected 5 VLA that support multi-step output and executed the 10 tasks in Section \ref{sec:task} on 30 types of distorted images. Note that among the five difficulty levels in Perception, we only executed the simplest one to ensure that the execution result of the reference image is correct. Thus, we ensured that the reason for execution failure came from the added distortion, not the image itself. Figure \ref{fig:real} shows examples of successful execution, results deviating from the ground truth, and emergency stops triggered by collisions with the table. We calculate the average Execution score under 30 distortion types and compare it with the Cognition and Decision scores, as shown in Figure \ref{fig:subject-bar} (b)(c), where findings are summarized as follows:

\CLA{Cognition} VS \CLA{Execution}: The SRCC of VLM results with the real world is less than 0.5. This corroborates the necessity of using VLA as subjects in the Embodied IQA task beyond VLM.
\\
\CLA{Decision} VS \CLA{Execution}: The SRCC of VLA results with the real world exceeds 0.6, indicating that Decision can represent Execution to some extent. However, this correlation is still not high enough, proving that certain real-world experiments are still indispensable for Embodied AI development. 
\\
\CLA{Perception}  VS \CLA{Cognition\&Decision}: Existing quality metrics have initially demonstrated the ability to handle Embodied IQA tasks, but there is still a gap compared to the traditional human-oriented paradigm. More advanced metrics should be developed in the upcoming Embodied AI era.

\begin{figure*}[t]
    \centering
    \subfigure[Cross Validation]{
    \begin{minipage}{.4\linewidth}
        \centering
        \vspace{-32mm}
        \renewcommand\arraystretch{1.25}
        \renewcommand\tabcolsep{2.7pt}
        \belowrulesep=0pt\aboverulesep=0pt
        \resizebox{\columnwidth}{!}{
        \begin{tabular}{l|cc|cc|cc}
    \toprule
    \multirow{2}{*}{Method} & \multicolumn{2}{c|}{LIVE \cite{dataset:live}} & \multicolumn{2}{c|}{TID \cite{dataset:tid}} & \multicolumn{2}{c}{VLM} \\
    \cline{2-7}
    & SRCC$\uparrow$ & PLCC$\uparrow$ & SRCC$\uparrow$ &  PLCC$\uparrow$ & SRCC$\uparrow$ & PLCC$\uparrow$ \\
    \midrule
    AHIQ \cite{quality:ahiq}     &0.5746&0.6402&0.2477&0.1147& 0.7240&0.6902\\
    CKDN \cite{quality:ckdn}     &0.5154&0.5379&0.2352&0.2316& 0.6977&0.6798\\
    DISTS \cite{quality:dists}   &0.7074&0.7075&0.7106&0.7570& 0.4307&0.2102\\
    LPIPS \cite{quality:lpips}   &0.5449&0.5713&0.3025&0.2708& 0.3430&0.3340\\
    TOPIQ-FR \cite{quality:topiq}&0.7285&0.7550&0.4510&0.2277& 0.7115&0.6788\\
    \cdashline{1-7}
    CLIPIQA \cite{quality:CLIPIQA}     &0.0472&0.0530&0.0222&0.0044& 0.0245&0.0307\\
    CNNIQA \cite{quality:CNNIQA}       &0.3477&0.4218&0.1074&0.0383& 0.4962&0.4755\\
    DBCNN \cite{quality:DBCNN}         &0.4801&0.5806&0.1891&0.0270& 0.5964&0.5855\\
    QualiClip \cite{quality:qualiclip} &0.6018&0.6531&0.2472&0.0769& 0.5890&0.5648\\
    TOPIQ-NR \cite{quality:topiq}      &0.5253&0.5577&0.2183&0.1291& 0.6806&0.6357\\
    \bottomrule
    \end{tabular}
        }
    \end{minipage}\hfill}\hskip.2em
    \subfigure[VLM \& Real-world]{\includegraphics[width=0.28\textwidth]{figs/Real-VLM.pdf}}\vspace{-3mm}
    \subfigure[VLA \& Real-world]{\includegraphics[width=0.28\textwidth]{figs/Real-VLA.pdf}}
    %\vspace{-4pt}
    \caption{Cross database validation on \CLB{Cognition} and human-oriented score, and the correlation between \CLB{VLM}\&\CLA{VLA} and Real-world. {\color[HTML]{FF0000} \textbf{\ding{72}}} denotes distortions with greater Real-world impact.}
    %\vspace{-2mm}
    \label{fig:subject-bar}
    \vspace{-3mm}
\end{figure*}

\begin{figure}[t]
    \centering
    \includegraphics[width=\linewidth]{figs/real-example.pdf}
    \vspace{-6mm}
    \caption{Positive/Negative cases of real-world experiment. The score of successful examples is 100, and deduct the Euclidean distance for failed examples. Triggering interruption will be scored as 0.}
    \label{fig:real}
    \vspace{-4mm}
\end{figure}

\section{Conclusion}

In this paper, we extend the application of IQA from a traditional human-oriented paradigm to Embodied AI.
To study which distortions have a negative impact on Embodied AI, we built a Perception-Cognition-Decision-Execution pipeline based on Mertonian Law and established a database for Embodied subjective preferences. Experiments show that more advanced IQA methods are needed to identify quality degradation for Embodied AI. We sincerely hope this Embodied IQA task can promote the application of Robotic Intelligence under complex distortions in the Real-world.

% \begin{table*}[t]
%     \centering
%     \renewcommand\arraystretch{1.25}
%     \renewcommand\tabcolsep{2.7pt}
%     \belowrulesep=0pt\aboverulesep=0pt
%     \caption{Cross-dataset Performance Comparison (SRCC, KRCC, PLCC)}
%     \label{tab:exp-content}
    
%     \resizebox{\linewidth}{!}{
%     \begin{tabular}{l|ccc|ccc|ccc}
%     \toprule
%     \multirow{2}{*}{Method} & \multicolumn{3}{c|}{VLM} & \multicolumn{3}{c|}{TID} & \multicolumn{3}{c}{LIVE} \\
%     \cline{2-10}
%     & SRCC$\uparrow$ & KRCC$\uparrow$ & PLCC$\uparrow$ & SRCC$\uparrow$ & KRCC$\uparrow$ & PLCC$\uparrow$ & SRCC$\uparrow$ & KRCC$\uparrow$ & PLCC$\uparrow$ \\
%     \midrule
%     AHIQ \cite{quality:ahiq}     & 0.7240&0.5172&0.6902&0.5746&0.4109&0.6402&0.2477&0.1835&0.1147
%  \\
%     CKDN \cite{quality:ckdn}     & 0.6977&0.4925&0.6798&0.5154&0.3653&0.5379&0.2352&0.1652&0.2316
%  \\
%     DISTS \cite{quality:dists}   & 0.7074&0.5182&0.7075&0.7106&0.5301&0.7570&0.4307&0.2991&0.2102
%  \\
%     LPIPS \cite{quality:lpips}   & 0.3430&0.2233&0.3340&0.0449&0.0363&0.0713&0.3025&0.2125&0.0708
%  \\
%     TOPIQ-FR \cite{quality:topiq}& 0.7115&0.5049&0.6788&0.7285&0.5391&0.7550&0.4510&0.3147&0.2277
%  \\
%     \cdashline{1-10}
%     CLIPIQA \cite{quality:CLIPIQA}     & 0.0245&0.0145&0.0307&0.0472&0.0320&0.0530&0.0222&0.0160&0.0044
%  \\
%     CNNIQA \cite{quality:CNNIQA}       & 0.4962&0.3413&0.4755&0.3477&0.2387&0.4218&0.1074&0.0819&0.0383
%  \\
%     DBCNN \cite{quality:DBCNN}         & 0.5964&0.4179&0.5855&0.4801&0.3393&0.5806&0.1891&0.1705&0.0270
%  \\
%     QualiClip \cite{quality:qualiclip} & 0.5890&0.4124&0.5648&0.6018&0.4282&0.6531&0.2472&0.1779&0.0769
%  \\
%     TOPIQ-NR \cite{quality:topiq}      & 0.6806&0.4770&0.6357&0.5253&0.3694&0.5577&0.2183&0.1555&0.1291
%  \\
%     \bottomrule
%     \end{tabular}}
    
%     \vspace{-4mm}
% \end{table*}





%%%%%%%%%%%%%%%%%%%%%%%%%%%%%%%%%%%%%%%%%%%%%%%%%%%%%%%%%%%%

\appendix

% \section{Technical Appendices and Supplementary Material}
% Technical appendices with additional results, figures, graphs and proofs may be submitted with the paper submission before the full submission deadline (see above), or as a separate PDF in the ZIP file below before the supplementary material deadline. There is no page limit for the technical appendices.

\section{Limitation \& Broader Impact}
% Chunyi
\textbf{[Limitation 1]}: As the first Embodied IQA work, we simplify the Perception to vision and Execution to robotic arms. Since the visual signals processed by humans account for 80\% of the total signals, while upper limb movements account for 50\% of all movements. Considering the consistency between humans and machines and the current limitations of Embodied AI, our simplification is reasonable. This will not affect the current main applications of Embodied AI, such as industrial assembly and home services. After the future vision-tactile fusion (Perception), quadruped robot dog (Execution), and other task scenarios are improved, we will further update the quality assessment data.

\textbf{[Limitation 2]}: The scale of our Real-world data is relatively small compared to Cognition and Decision. In the VLM and VLA steps, we already have the largest amount of data (millions) compared to previous IQA work. This is because the high cost of real machine data requires a lot of manpower in the site layout and verification stages. Although 1,500 annotation samples are not enough as quality labels, they are sufficient to verify the Sim2Real consistency of Cognition-Execution and Decision-Execution. With the further development of Embodied AI, we believe an automated Real-world pipeline will be developed, from which we will expand the scale of Execution labels.

\textbf{[Broader Impact] (Positive)}: IQA can expand the application scenarios of Embodied AI, extending it from the in-lab environment to distortions in the Real-world. We collect the subjective preferences of Embodied AI, thus objectively judge the `utility' of images before executing specific tasks. In this way, distorted images such as jitter and blur can be effectively filtered. Such quality indicators can be used for all visual applications for Embodied AI, such as video coding, super-resolution, defogging, denoising, etc. Considering that the amount of visual signals consumed by machines has exceeded humans since 2023, visual quality indicators for Embodied AI can fill this research gap.

\textbf{[Broader Impact] (Negative)}: The general use of visual quality indicators in Embodied AI may affect traditional human-oriented tasks. Considering that humans, VLM (Cognition), and VLA (Decision) have different preferences, only evaluating the preferences of VLM and VLA will inevitably lead to scores that are not relevant to humans. Therefore, in future international protocols, it is recommended to integrate the three IQA paradigms for humans, VLM, and VLA together, and select appropriate quality indicators based on the user end.

% 正面影响在扩展Embodied AI的应用场景，将其从in-lab的环境延申到真实世界的失真中。我们总结了Embodied AI的主观偏好，在执行具体任务前，就可以判断图像的utility。由此对抖动、模糊等失真图像，实现有效的过滤。这样的质量指标可以用于一切面向Embodied AI的视觉应用，例如视频编码通信，超分，去雾，去噪等。考虑到自2023年起，机器消费的视觉信号量已超过人类，面向Embodied AI的视觉质量指标可以弥补这一研究空白。
% 负面影响在Embodied AI的视觉质量指标的泛用，可能影响传统的human-oriented task。考虑到human，VLM（Cognition），VLA (Decision)的偏好各不相同，仅评估VLM与VLA的偏好，势必导致分数与人类不够相关。因此在未来的标准中，建议将面向human，VLM，VLA的三种IQA范式同时考量，根据用户端选择合适的质量指标。


\section{Robotics Settings}
% Jianbo
This section provides a detailed derivation of the forward and inverse kinematics for the Universal Robots (UR5), a 6-DoF collaborative robot. The Denavit-Hartenberg (D-H) convention is used to establish the kinematic model. 

In the experiments part, the initial pose is obtained through forward kinematics by recording the initial rotational angles of the six joints and calculating the end-effector's pose relative to the base coordinate frame. The incremental pose output by the VLA is then multiplied with the initial pose to derive the step-by-step poses. The robotic arm's actual motion is resolved via inverse kinematics, which computes the required rotational angles for each joint motor to achieve the target configuration.

The D-H parameters define the geometry of the robot manipulator by establishing a coordinate frame $\{i\}$ attached to each link $i$. The transformation from frame $\{i-1\}$ to frame $\{i\}$, denoted $A_i^{i-1}$, is described by four parameters associated with link $i-1$ and joint $i$:
\begin{itemize}
    \item $\theta_i$: Joint Angle - the rotation about the $z_{i-1}$ axis, from $x_{i-1}$ to $x_i$. For a revolute joint, $\theta_i$ is the joint variable.
    \item $d_i$: Link Offset - the distance along the $z_{i-1}$ axis from the origin of frame $\{i-1\}$ to the intersection of the $z_{i-1}$ axis with the $x_i$ axis. For a prismatic joint, $d_i$ is the joint variable.
    \item $a_{i-1}$: Link Length - the distance along the $x_i$ axis (which is the common normal between $z_{i-1}$ and $z_i$) from the intersection of $z_{i-1}$ and $x_i$ axis to the origin of frame $\{i\}$.
    \item $\alpha_{i-1}$: Link Twist - the angle about the $x_i$ axis, from $z_{i-1}$ to $z_i$.
\end{itemize}

The UR5 D-H parameters used in this paper shown in Table \ref{tab:dh_params}. Where $a_2, a_3$ are physical link lengths associated with links 2 and 3 respectively (used as $a_{i-1}$ parameters in the table for joints 3 and 4), and $d_1, d_4, d_5, d_6$ are link offsets. The $\theta_i^*$ are the joint variables.

\begin{table}[t]
    \centering
    \renewcommand\arraystretch{1.25}
    \renewcommand\tabcolsep{7pt}
    \belowrulesep=0pt\aboverulesep=0pt
    \caption{Parameter settings of Robotic arm UR5 D-H. The specific action depends on 6 frames.}
    \label{tab:dh_params}
    \resizebox{0.6\linewidth}{!}{
    \begin{tabular}{ccccc}
        \toprule
        Joint Frame $i$ & $\alpha_{i-1}$ (rad) & $a_{i-1}$ (m) & $d_i$ (m) & $\theta_i$ (rad) \\
        \midrule
        1         & $0$                  & $0$           & $d_1$     & $\theta_1^*$     \\
        2         & $\pi/2$              & $0$           & $0$       & $\theta_2^*$     \\
        3         & $0$                  & $a_2$     &     $0$    &    $\theta_3^*$     \\
        4         & $0$      &           $a_3$     &     $d_4$    &  $\theta_4^*$     \\
        5         & $\pi/2$           &    $0$          &  $d_5$     & $\theta_5^*$     \\
        6         & $-\pi/2$           &   $0$          &  $d_6$    &  $\theta_6^*$     \\
        \bottomrule
    \end{tabular}}
\end{table}

Typical UR5 parameter values (example, signs depend on coordinate frame choices):
$d_1 = 0.089159$ m, $a_2 = 0.42500$ m (often negative in some tables: $-0.42500$), $a_3 = 0.39225$ m (often negative: $-0.39225$), $d_4 = 0.10915$ m, $d_5 = 0.09465$ m, $d_6 = 0.0823$ m.

\section{Forward Kinematics}
The standard D-H transformation matrix $A_i^{i-1}$ from frame $\{i-1\}$ to frame $\{i\}$ is defined as a product of four basic transformations:
\begin{equation}
A_i^{i-1} = \RotZ{\theta_i} \TransZ{d_i} \TransX{a_{i-1}} \RotX{\alpha_{i-1}},
\end{equation}
\begin{equation}
A_i^{i-1} = \mat{\ct{i} & -\st{i}\ca{i-1} & \st{i}\sa{i-1} & a_{i-1}\ct{i} \\
                  \st{i} & \ct{i}\ca{i-1} & -\ct{i}\sa{i-1} & a_{i-1}\st{i} \\
                  0       & \sa{i-1}        & \ca{i-1}        & d_i \\
                  0       & 0               & 0               & 1 },
\end{equation}                  

where $\mathbf{Tr}(\cdot)$ and $\mathbf{R}(\cdot)$ denote the trajectory and rotation matrix projected on a certain axis.
For simplicity, the following symbols will be defined in the subsequent sections: $c_i = \cos(\theta_i)$, $s_i = \sin(\theta_i)$. $c_{ij} = \cos(\theta_i+\theta_j)$, $s_{ij} = \sin(\theta_i+\theta_j)$.
Using the D-H parameters from Table \ref{tab:dh_params}, the individual transformation matrix $A_i^{i-1}$ for robotic manipulators are:

    \begin{equation}
    A_1^0 = \mat{\ct{1} & -\st{1} & 0 & 0 \\ \st{1} & \ct{1} & 0 & 0 \\ 0 & 0 & 1 & d_1 \\ 0 & 0 & 0 & 1},
    \end{equation}
    

    \begin{equation}
    A_2^1 = \mat{\ct{2} & 0 & \st{2} & 0 \\ \st{2} & 0 & -\ct{2} & 0 \\ 0 & 1 & 0 & 0 \\ 0 & 0 & 0 & 1},
    \end{equation}
    

    \begin{equation}
    A_3^2 = \mat{\ct{3} & -\st{3} & 0 & a_2\ct{3} \\ \st{3} & \ct{3} & 0 & a_2\st{3} \\ 0 & 0 & 1 & 0 \\ 0 & 0 & 0 & 1},
    \end{equation}


    \begin{equation}
    A_4^3 = \mat{\ct{4} & -\st{4} & 0 & a_3\ct{4} \\ \st{4} & \ct{4} & 0 & a_3\st{4} \\ 0 & 0 & 1 & d_4 \\ 0 & 0 & 0 & 1},
    \end{equation}


    \begin{equation}
    A_5^4 = \mat{\ct{5} & 0 & \st{5} & 0 \\ \st{5} & 0 & -\ct{5} & 0 \\ 0 & 1 & 0 & d_5 \\ 0 & 0 & 0 & 1},
    \end{equation}


    \begin{equation}
    A_6^5 = \mat{\ct{6} & 0 & -\st{6} & 0 \\ \st{6} & 0 & \ct{6} & 0 \\ 0 & -1 & 0 & d_6 \\ 0 & 0 & 0 & 1}.
    \end{equation}

The total transformation matrix $T_0^6$ from the base frame $\{0\}$ to the end-effector frame $\{6\}$ is:
\begin{equation}
T_0^6 = A_1^0 A_2^1 A_3^2 A_4^3 A_5^4 A_6^5 = \mat{R_0^6 & p_0^6 \\ \mathbf{0}_{1 \times 3} & 1} = \mat{n_x & s_x & a_x & p_x \\ n_y & s_y & a_y & p_y \\ n_z & s_z & a_z & p_z \\ 0 & 0 & 0 & 1},
\end{equation}
where $R_0^6 = [n, s, a]$ refers to the rotation matrix part of $T_0^6$, where $n=[n_x,n_y,n_z]\T$, $s=[s_x,s_y,s_z]\T$, and $a=[a_x,a_y,a_z]\T$ are column vectors representing the $x, y, z$ axes of frame $\{6\}$ expressed in frame $\{0\}$, respectively.
$p_0^6 = [p_x, p_y, p_z]\T$ represents the translation vector part of $T_0^6$, namely the position of the origin of frame $\{6\}$ expressed in frame $\{0\}$.

\section{Inverse Kinematics}
The objective of Inverse Kinematics (IK) is to determine the set of joint angles $(\theta_1, \dots, \theta_6)$ that achieve a desired end-effector pose $T_0^6$. The UR5 possesses a spherical wrist (axes of joints 4, 5, and 6 intersect at a common point, the wrist center), which allows for a decoupled analytical solution. First, the position of the wrist center is determined, which allows solving for the first three joints. Then, the orientation of the end-effector is used to solve for the remaining three wrist joints.

\subsection{Calculation of the Wrist Center Point}
The wrist center point (WCP), $p_{wc}$, is typically defined as the origin of frame $\{5\}$. Its position can be found by translating from the end-effector origin, $p_0^6$, backwards along the approach vector $a$ (the $z_6$-axis expressed in frame $\{0\}$) by a distance $d_6$:
\begin{equation}
p_{wc} = p_0^6 - R_0^6 \mat{0 \\ 0 \\ d_6} = \mat{p_x \\ p_y \\ p_z} - d_6 \mat{a_x \\ a_y \\ a_z} = \mat{p_x - d_6 a_x \\ p_y - d_6 a_y \\ p_z - d_6 a_z},
\end{equation}
where $p_{wc} = (x_{wc}, y_{wc}, z_{wc})\T$ means position vector of the wrist center point in frame $\{0\}$.

\subsection{Solving for Base Joints}
The $y$-coordinate of $p_{wc}$ when expressed in frame $\{1\}$, denoted $p_{wc,y}^1$, can be shown to be $d_4+d_5$ for this specific D-H parameterization (where $p_{wc}$ is the origin of frame $\{5\}$).
We have $p_{wc,y}^1 = -x_{wc} \st{1} + y_{wc} \ct{1} = d_4+d_5$.
This equation can be solved for $\theta_1$:
\begin{equation}
\theta_1 = \atan{y_{wc}}{x_{wc}} - \atan{d_4+d_5}{\sigma_1 \sqrt{x_{wc}^2 + y_{wc}^2 - (d_4+d_5)^2}},
\end{equation}
where $\sigma_1 = \pm 1$ denotes two possible solutions for $\theta_1$. 
Function $\rm{atan2}(\cdot,\cdot)$ transforms
two Cartesian coordinates
to polar.
If the term under the square root is negative, the target $p_{wc}$ is unreachable.

\subsection{Solving for Elbow Joints}
With $\theta_1$ known, transform $p_{wc}$ into frame $\{1\}$. Let $K_x$ and $K_z$ be coordinates of $p_{wc}$ relevant for the planar geometry of links 2 and 3:
$K_x = x_{wc} \ct{1} + y_{wc} \st{1}$
$K_z = z_{wc} - d_1$. From the geometry of the first three links (considering $\alpha_1=\pi/2$ which introduces a rotation making $z_1$ horizontal in the new projected plane if $\theta_2=0$):
$K_x = a_2 \ct{2} + a_3 \ctt{2}{3}$
$K_z = -a_2 \st{2} - a_3 \stt{2}{3}$. Squaring and adding these two equations yields:
$K_x^2 + K_z^2 = a_2^2 + a_3^2 + 2 a_2 a_3 \ct{3}$. This allows solving for $\theta_3$:
\begin{equation}
\ct{3} = \frac{K_x^2 + K_z^2 - a_2^2 - a_3^2}{2 a_2 a_3},
\end{equation}
where $\sigma_3 = \pm 1$ corresponds to up/down configurations. From $\st{3} = \sigma_3 \sqrt{1-\ct{3}^2}$ we have:
\begin{equation}
    \theta_3 = \atan{\st{3}}{\ct{3}}.
\end{equation}
To solve for $\theta_2$, rearrange the equations for $K_x$ and $K_z$:
Let $k_1 = a_2 + a_3\ct{3}$ and $k_2 = a_3\st{3}$.
Then $K_x = k_1 \ct{2} - k_2 \st{2}$ and $K_z = -k_1 \st{2} - k_2 \ct{2}$. 
Solving this system for $\st{2}$ and $\ct{2}$:
$\st{2} = -(k_1 K_z + k_2 K_x) / (k_1^2+k_2^2)$
$\ct{2} = (k_1 K_x - k_2 K_z) / (k_1^2+k_2^2)$ (Note: $k_1^2+k_2^2 = K_x^2+K_z^2$):
\begin{equation}
    \theta_2 = \atan{-(k_1 K_z + k_2 K_x)}{k_1 K_x - k_2 K_z}.
\end{equation}
Alternatively, a more robust form is often $\theta_2 = \atan{-K_z}{K_x} - \atan{k_2}{k_1}$.

\subsection{Solving for Wrist Joints}
Once $\theta_1, \theta_2, \theta_3$ are known, the rotation matrix $R_0^3$ from the base to frame $\{3\}$ can be computed: $R_0^3 = (A_1^0 A_2^1 A_3^2)_{rot}$.
The rotation matrix from frame $\{3\}$ to frame $\{6\}$ is then $R_3^6 = (R_0^3)\T R_0^6$. Let $R_3^6 = [r'_{ij}]$.
The matrix $R_3^6$ can also be expressed as the product of rotations for joints 4, 5, 6 using their D-H parameters:
$R_3^6 = \RotZ{\theta_4}\RotX{\alpha_3} \RotZ{\theta_5}\RotX{\alpha_4} \RotZ{\theta_6}\RotX{\alpha_5}$.
For the UR5 D-H parameters in Table \ref{tab:dh_params}:
$R_3^6 = \RotZ{\theta_4} \RotZ{\theta_5}\RotX{\pi/2} \RotZ{\theta_6}\RotX{-\pi/2}$.
The symbolic product is:
\begin{equation}
R_3^6 = \mat{ c_4c_5c_6-s_4s_5c_6 & c_4s_5+s_4c_5 & c_4c_5s_6-s_4s_5s_6 \\ s_4c_5c_6+c_4s_5c_6 & s_4s_5-c_4c_5 & s_4c_5s_6+c_4s_5s_6 \\ s_6 & 0 & -c_6 }.
\end{equation}
By comparing elements of the numerically computed $R_3^6 = [r'_{ij}]$ with this symbolic form:
\begin{enumerate}
    \item $r'_{32}$ must be $0$. If the computed $( (R_0^3)\T R_0^6 )_{32}$ is significantly non-zero, it indicates no solution for this wrist structure or a modeling error.
    \item From $r'_{31} = s_6$ and $r'_{33} = -c_6$:
    \begin{equation}
    \theta_6 = \atan{r'_{31}}{-r'_{33}}.
    \end{equation}
    This provides a unique solution for $\theta_6$ in $(-\pi, \pi]$. Another solution is $\theta_6 \pm \pi$ (if $s_6, c_6$ are flipped), but usually we seek solutions within joint limits.
    \item From $r'_{12} = c_4s_5+s_4c_5 = s_{4+5}$ and $r'_{22} = s_4s_5-c_4c_5 = -c_{4+5}$:
    \begin{equation}
    \theta_4+\theta_5 = \atan{r'_{12}}{-r'_{22}}.
    \end{equation}
    Consider the movement of these two elements as a whole, we have $\phi_{45} = \theta_4+\theta_5$.
    \item To find $\theta_5$ and $\theta_4$ separately:
    $s_5 = c_4 r'_{12} + s_4 r'_{22} = c_4 s_{\phi_{45}} + s_4 (-c_{\phi_{45}}) = \sin(\phi_{45}-\theta_4) = \sin\theta_5$.
    $c_5 = s_4 r'_{12} - c_4 r'_{22} = s_4 s_{\phi_{45}} - c_4 (-c_{\phi_{45}}) = \cos(\phi_{45}-\theta_4) = \cos\theta_5$.
    A common method to solve for $\theta_5$ (wrist roll) for many spherical wrists involves:
    \begin{equation}
        \theta_5 = \sigma_5 \arccos\left( \frac{r'_{11} + r'_{22} \frac{s_4}{c_4} }{c_6(c_4 - s_4 \frac{s_4}{c_4})} \right).
    \end{equation}
    According to $\theta_5$ we have $\theta_4=\phi_{45}-\theta_5$. Thus all rotation angles can be retrieved.
\end{enumerate}
Typically, UR5 has 8 unique inverse kinematics solutions ($\sigma_1=\pm 1, \sigma_3=\pm 1, \sigma_5=\pm 1$ for the choice of $s_5$). Singularities (e.g., $s_5=0$) lead to infinite solutions where $\theta_4$ and $\theta_6$ are coupled.

\subsection{Handling Singularities}
\begin{itemize}
    \item Shoulder Singularity: Occurs if $x_{wc}^2 + y_{wc}^2 - (d_4+d_5)^2 = 0$. The wrist center lies on the $z_0$ axis (for $d_4+d_5=0$) or a cylinder around $z_0$. $\theta_1$ is not uniquely defined.
    \item Elbow Singularity: Occurs if $K_x^2 + K_z^2 - a_2^2 - a_3^2 = \pm 2 a_2 a_3$, meaning $\ct{3}=\pm 1$ (arm fully extended or folded). $\st{3}=0$, so $\theta_2$ solution becomes simpler but an infinite number of $\theta_2$ might exist if $p_{wc}$ is on $z_1$.
    \item Wrist Singularity: Occurs if $s_5=0$ (i.e., $\theta_5=0$ or $\pi$). Axes $z_4$ and $z_6$ align. In this case, only the sum or difference $(\theta_4 \pm \theta_6)$ can be determined according to Section D.4. One angle can be chosen arbitrarily, and the other is then fixed.
\end{itemize}

\begin{figure}[t]
    \centering
    \subfigure[Annotation]{\includegraphics[width=0.48\textwidth]{supp/interface-crop-1.png}}\vspace{-3mm}
    \subfigure[Verification]{\includegraphics[width=0.48\textwidth]{supp/interface-crop-2.png}}
    \caption{Human annotation and verification Interface. Subjects will cyclically define five tasks with different contents and increasing difficulty, and submit them to robotic experts for verification.}
    \label{fig:interface}
    \vspace{-4mm}
\end{figure}

\section{Subjective Perception Task Definition}
% Jiahao，界面
Before VLM and VLA inference, we organized five Ph. D. candidates as a panel to define five downstream tasks for each image as shown in Figure \ref{fig:interface}.
To avoid bias from a single subject, each sample is sent to five subjects in a random order to design specific tasks based on the image. 
The 1,230 samples to be annotated come from seven Robotic database \cite{dataset:robot-droid,dataset:robot-kuka,dataset:robot-ssvtp,dataset:robot-Jacquard,dataset:robot-ocid_vlg,dataset:robot-maniskill,dataset:robot-taco_play1}, and are filtered according to the settings in the main text, to retain only high-quality images.
Subjects can see the previous tasks, and the new tasks they design need to be more difficult than them and test different abilities (for example, if an object has been pushed, try not to grab it again). Images with display errors, NSFW, or offensive content will be removed. After each image has five task labels, a professional robotics engineer will adjust the specific samples. Based on operational experience, the difficulty of the five tasks will be re-ranked, and unreasonable tasks will be modified.

\begin{figure*}[t]
    \centering
    \begin{minipage}[]{0.21\linewidth}
    \centering
    \centerline{\includegraphics[width = \textwidth]{figs/DimVLM.pdf}}
    % \centerline{VQA}\medskip
    \end{minipage}
    \begin{minipage}[]{0.75\linewidth}
    \centering
    \centerline{\includegraphics[width = \textwidth]{figs/vlm-subtypes.pdf}}
    % \centerline{CAP}\medskip
    \end{minipage}
    \vspace{-2mm}

    \caption{Correlation between the general Cognition score and the 3 dimensions from VLMs, and the score distribution in Sim2Real, First/Third perspectives, Main object, and Background sub-categories.}
    \label{fig:violin-vlm}
    \vspace{-3mm}
\end{figure*}

\section{Cognition IQA Experiment}
% Jiahao，VLM反过来的实验
Due to space limitations in the main text, we mainly discuss the Decision step (specific to Embodied AI), and the Cognition step (common to general machines) is listed in this section. First, the Cognition score given by VLM is shown in Figure \ref{fig:violin-vlm}. Compared with the three dimensions of Decision in Figure \ref{fig:violin}, the correlation between the Cognition dimensions is higher, and the distribution difference between different categories of data is smaller. This fully demonstrates the difference between the reasoning mechanisms of VLM and VLA, and proves the rationality of separating these two steps.

Therefore, in addition to Decision, we also conducted IQA experiments on Cognition, following the training/testing settings in the main text. Table \ref{tab:exp-content-2} presents the performance of advanced quality metrics on Cognition, compared with Decision in Table \ref{tab:exp-content}, current IQA metrics has better prediction results on Cognition. Since the current IQA method is more related to VLM than VLA, the quality indicators that general machines already have are initially available, but Embodied AI cannot be effectively evaluated. It is worth mentioning that the zero-shot baseline method on Cognition can occasionally even achieve an SRCC of more than 0.7, surpassing a number of fine-tuned methods; while the baseline on Decision is significantly weaker. This is exactly why we separated the Robot Visual System from the Machine Visual System and used the Mortonian system to model the Intelligent System in four steps. In short, we hope that the Embodied IQA database can promote more complete quality indicators, whether applied for VLM or VLA as subjects in Embodied tasks.


\begin{table*}[t]
    \centering
    \renewcommand\arraystretch{1.25}
    \renewcommand\tabcolsep{2.7pt}
    \belowrulesep=0pt\aboverulesep=0pt
    \caption{Using 15 advanced IQA metrics to evaluate the \CLB{Cognition} score from VLMs, including zero-shot, FR, and NR metrics. [Keys: \CLA{Best}/\CLB{Second best} in group; \textbf{Baseline}; \colorbox{gray!20}{Lower} than baseline.]}
    \label{tab:exp-content-2}

    \resizebox{\linewidth}{!}{
\begin{tabular}{l|l|ccc|ccc|ccc|ccc|ccc}
\toprule
\multicolumn{2}{c|}{Dimension}            & \multicolumn{3}{c|}{Precision} & \multicolumn{3}{c|}{Recall}  & \multicolumn{3}{c|}{Semantic}  & \multicolumn{3}{c|}{First Perspective}& \multicolumn{3}{c}{Third Perspective}   \\ \cline{1-17}
Group                 & \multicolumn{1}{c|}{Metric}        & SRCC$\uparrow$   & KRCC$\uparrow$   & PLCC$\uparrow$   & SRCC$\uparrow$   & KRCC$\uparrow$   & PLCC$\uparrow$   & SRCC$\uparrow$   & KRCC$\uparrow$   & PLCC$\uparrow$   & SRCC$\uparrow$   & KRCC$\uparrow$   & PLCC$\uparrow$   & SRCC$\uparrow$   & KRCC$\uparrow$   & PLCC$\uparrow$   \\ \midrule
\multirow{5}{*}{Zero} & PSNR          &0.3432&0.2339&0.3661&0.3186&0.2168&0.3369&0.3257&0.2225&0.3520&0.4142&0.2831&0.4393&0.3605&0.2477&0.4140


 \\
                      & SSIM \cite{quality:ssim}         & 0.5809&0.4055&0.5561&0.5558&0.3871&0.5241&0.5798&0.4057&0.5521&0.6244&0.4383&0.5805&0.5849&0.4094&0.5438


 \\
                      & Brisque \cite{quality:brisque}      & 0.3527&0.2380&0.3342&0.3537&0.2412&0.3319&0.3596&0.2442&0.3375&0.3232&0.2187&0.2987&0.3751&0.2545&0.3554


 \\
                      & Q-Align \cite{quality:q-align}  & \textbf{0.7045}&\textbf{0.5067}&\textbf{0.6721}&\textbf{0.6755}&\textbf{0.4798}&\textbf{0.6278}&\textbf{0.7040}&\textbf{0.5058}&\textbf{0.6687}&\textbf{0.6622}&\textbf{0.4767}&\textbf{0.6214}&\textbf{0.7321}&\textbf{0.5349}&\textbf{0.6970}

\\
                      & Q-Align+ \cite{quality:q-align+} & 0.4697&0.3184&0.4049&0.4524&0.3063&0.3623&0.4722&0.3202&0.3958&0.4687&0.3188&0.4263&0.5082&0.3421&0.4351


 \\
\cdashline{1-17}
\multirow{5}{*}{FR}   & AHIQ \cite{quality:ahiq}         & \CLB{0.7941} & \CLB{0.5930} & \CLB{0.7901} & \CLB{0.7747} & \CLB{0.5753} & \CLB{0.7683} & \CLB{0.7983} & \CLB{0.5976} & \CLB{0.7919} & \CLA{0.8035} & \CLA{0.6039} & \CLA{0.8015} & \CLB{0.8288} & \CLB{0.6342} & \CLB{0.8292}


\\
                      & CKDN \cite{quality:ckdn}         & 0.7461&0.5460&0.7444&0.7387&0.5380&0.7332&0.7516&0.5508&0.7470&0.7556&0.5587&0.7547&0.7836&0.5808&0.7810

\\
                      & DISTS \cite{quality:dists}        & \colorbox{gray!20}{0.7017}&0.5128&0.7052&0.6887&0.5010&0.6863&0.7080&0.5188&0.7096&0.7307&0.5424&0.7299&0.7535&0.5584&0.7530

 \\
                      & LPIPS \cite{quality:lpips}        & \colorbox{gray!20}{0.6681}&\colorbox{gray!20}{0.4785}&\colorbox{gray!20}{0.6165}&\colorbox{gray!20}{0.6463}&\colorbox{gray!20}{0.4610}&\colorbox{gray!20}{0.5975}&\colorbox{gray!20}{0.6714}&\colorbox{gray!20}{0.4812}&\colorbox{gray!20}{0.6179}&0.6797&0.4929&0.6292&\colorbox{gray!20}{0.6893}&\colorbox{gray!20}{0.5008}&\colorbox{gray!20}{0.6485}

 \\
                      & TOPIQ-FR \cite{quality:topiq}     & \CLA{0.8209}&\CLA{0.6241}&\CLA{0.8194}&\CLA{0.8160}&\CLA{0.6170}&\CLA{0.8126}&\CLA{0.8326}&\CLA{0.6363}&\CLA{0.8289}&\CLB{0.7997}&\CLB{0.5967}&\CLB{0.7932}&\CLA{0.8521}&\CLA{0.6574}&\CLA{0.8462}

\\
\cdashline{1-17}
\multirow{6}{*}{NR}   & CLIPIQA \cite{quality:CLIPIQA}      & \colorbox{gray!20}{0.3111}&\colorbox{gray!20}{0.2101}&\colorbox{gray!20}{0.3185}&\colorbox{gray!20}{0.3013}&\colorbox{gray!20}{0.2038}&\colorbox{gray!20}{0.3141}&\colorbox{gray!20}{0.3167}&\colorbox{gray!20}{0.2156}&\colorbox{gray!20}{0.3257}&\colorbox{gray!20}{0.1748}&\colorbox{gray!20}{0.1198}&\colorbox{gray!20}{0.1778}&\colorbox{gray!20}{0.3889}&\colorbox{gray!20}{0.2620}&\colorbox{gray!20}{0.3821}



 \\
                      & CNNIQA \cite{quality:CNNIQA}       & \colorbox{gray!20}{0.4864}&\colorbox{gray!20}{0.3359}&\colorbox{gray!20}{0.4818}&\colorbox{gray!20}{0.4793}&\colorbox{gray!20}{0.3312}&\colorbox{gray!20}{0.4761}&\colorbox{gray!20}{0.4880}&\colorbox{gray!20}{0.3368}&\colorbox{gray!20}{0.4861}&\colorbox{gray!20}{0.4719}&\colorbox{gray!20}{0.3247}&\colorbox{gray!20}{0.4735}&\colorbox{gray!20}{0.5336}&\colorbox{gray!20}{0.3744}&\colorbox{gray!20}{0.5431}



\\
                      & DBCNN \cite{quality:DBCNN}        & \colorbox{gray!20}{0.5687}&\colorbox{gray!20}{0.3921}&\colorbox{gray!20}{0.5553}&\colorbox{gray!20}{0.5349}&\colorbox{gray!20}{0.3650}&\colorbox{gray!20}{0.5183}&\colorbox{gray!20}{0.5596}&\colorbox{gray!20}{0.3835}&\colorbox{gray!20}{0.5453}&\colorbox{gray!20}{0.5341}&\colorbox{gray!20}{0.3578}&\colorbox{gray!20}{0.5346}&\colorbox{gray!20}{0.6622}&\colorbox{gray!20}{0.4699}&\colorbox{gray!20}{0.6489}



 \\
                      & QualiClip\cite{quality:qualiclip}     & \CLB{0.7416} & \CLB{0.5425} & \CLB{0.7389} & \CLB{0.7399} & \CLB{0.5399} & \CLB{0.7383} & \CLB{0.7524} & \CLB{0.5514} & \CLB{0.7499} & \CLB{0.6960} & \CLB{0.5022} & \CLB{0.6913} & \CLB{0.7864} & \CLB{0.5832} & \CLB{0.7711}



\\
                      & TOPIQ-NR \cite{quality:topiq}     & \CLA{0.7941}&\CLA{0.5933}&\CLA{0.7897}&\CLA{0.7818}&\CLA{0.5812}&\CLA{0.7761}&\CLA{0.8031}&\CLA{0.6039}&\CLA{0.7958}&\CLA{0.7854}&\CLA{0.5846}&\CLA{0.7819}&\CLA{0.8512}&\CLA{0.6577}&\CLA{0.8449}



 \\ \bottomrule
\end{tabular}}

\resizebox{\linewidth}{!}{
\begin{tabular}{l|l|ccc|ccc|ccc|ccc|ccc}
\toprule
\multicolumn{2}{c|}{Dimension}          & \multicolumn{3}{c|}{Mild Distortion}  & \multicolumn{3}{c|}{Medium Distortion} & \multicolumn{3}{c|}{Severe Distortion}  & \multicolumn{3}{c|}{Real-world} & \multicolumn{3}{c}{Simulation}   \\ \cline{1-17}
Group                 & \multicolumn{1}{c|}{Metric}        & SRCC$\uparrow$   & KRCC$\uparrow$   & PLCC$\uparrow$   & SRCC$\uparrow$   & KRCC$\uparrow$   & PLCC$\uparrow$   & SRCC$\uparrow$   & KRCC$\uparrow$   & PLCC$\uparrow$   & SRCC$\uparrow$   & KRCC$\uparrow$   & PLCC$\uparrow$   & SRCC$\uparrow$   & KRCC$\uparrow$   & PLCC$\uparrow$   \\ \midrule
\multirow{5}{*}{Zero} & PSNR          & 0.2350&0.1588&0.3725&0.1122&0.0751&0.1527&0.0811&0.0566&0.0763&0.4390&0.3056&0.4575&0.3613&0.2466&0.3945



 \\
                      & SSIM \cite{quality:ssim}         & 0.3263&0.2229&0.3133&0.1434&0.0972&0.1080&0.2102&0.1375&0.2402&0.6319&0.4507&0.5916&\textbf{0.6645}&\textbf{0.4688}&\textbf{0.6449}



 \\

                      & Brisque \cite{quality:brisque}      & 0.1525&0.1205&0.1017&0.1503&0.1168&0.0967&0.2088&0.1502&0.1334&0.4143&0.3907&0.2798&0.0636&0.0423&0.0421

 \\
                      & Q-Align \cite{quality:q-align} & \textbf{0.5092}&\textbf{0.3572}&\textbf{0.4793}&\textbf{0.4491}&\textbf{0.3071}&\textbf{0.4360}&\textbf{0.4056}&\textbf{0.2775}&\textbf{0.5047}&\textbf{0.7764}&\textbf{0.5739}&\textbf{0.7475}&0.6642&0.4648&0.6319

 \\
                      & Q-Align+ \cite{quality:q-align+} & 0.2995&0.2006&0.2657&0.3422&0.2268&0.2990&0.2154&0.1440&0.2240&0.5643&0.3898&0.5316&0.4804&0.3302&0.4709


 \\ \cdashline{1-17}
\multirow{5}{*}{FR}   & AHIQ \cite{quality:ahiq}         & \CLB{0.5921} & \CLB{0.4191} & \CLB{0.5842} & 0.5538 & \CLB{0.3895} & \CLB{0.5644} & \CLB{0.5580} & \CLB{0.3968} & \CLB{0.6064} & \CLB{0.8293} & \CLB{0.6339} & \CLB{0.8214} & \CLB{0.8138} & \CLB{0.6117} & \CLB{0.8041}



 \\
                      & CKDN \cite{quality:ckdn}         & 0.5497&0.3846&0.5672&\CLB{0.5539}&0.3872&0.5537&0.5087&0.3533&0.5287&\colorbox{gray!20}{0.7480}&\colorbox{gray!20}{0.5494}&\colorbox{gray!20}{0.7457}&0.7706&0.5651&0.7727


\\
                      & DISTS \cite{quality:dists}        & 0.5376&0.3796&0.5179&\colorbox{gray!20}{0.2957}&\colorbox{gray!20}{0.2030}&\colorbox{gray!20}{0.2894}&\colorbox{gray!20}{0.3258}&\colorbox{gray!20}{0.2215}&\colorbox{gray!20}{0.3657}&\colorbox{gray!20}{0.7606}&\colorbox{gray!20}{0.5695}&0.7566&0.7675&0.5686&0.7677


\\
                      & LPIPS \cite{quality:lpips}        & \colorbox{gray!20}{0.4635}&\colorbox{gray!20}{0.3229}&\colorbox{gray!20}{0.4285}&\colorbox{gray!20}{0.3876}&\colorbox{gray!20}{0.2661}&\colorbox{gray!20}{0.3700}&\colorbox{gray!20}{0.3168}&\colorbox{gray!20}{0.2148}&\colorbox{gray!20}{0.3225}&\colorbox{gray!20}{0.7128}&\colorbox{gray!20}{0.5190}&\colorbox{gray!20}{0.6493}&0.6693&0.4825&\colorbox{gray!20}{0.6030}



 \\
                      & TOPIQ-FR \cite{quality:topiq}     & \CLA{0.6755}&\CLA{0.4873}&\CLA{0.6661}&\CLA{0.6152}&\CLA{0.4355}&\CLA{0.6194}&\CLA{0.5798}&\CLA{0.4070}&\CLA{0.6125}&\CLA{0.8392}&\CLA{0.6398}&\CLA{0.8232}&\CLA{0.8449}&\CLA{0.6501}&\CLA{0.8442}


 \\ \cdashline{1-17}
\multirow{6}{*}{NR}   & CLIPIQA \cite{quality:CLIPIQA}      & \colorbox{gray!20}{0.1542}&\colorbox{gray!20}{0.1031}&\colorbox{gray!20}{0.1350}&\colorbox{gray!20}{0.0375}&\colorbox{gray!20}{0.0262}&\colorbox{gray!20}{0.0319}&\colorbox{gray!20}{0.0974}&\colorbox{gray!20}{0.0645}&\colorbox{gray!20}{0.1245}&\colorbox{gray!20}{0.3944}&\colorbox{gray!20}{0.2624}&\colorbox{gray!20}{0.3853}&\colorbox{gray!20}{0.6178}&\colorbox{gray!20}{0.4305}&\colorbox{gray!20}{0.5342}



 \\
                      & CNNIQA \cite{quality:CNNIQA}       & \colorbox{gray!20}{0.2697}&\colorbox{gray!20}{0.1797}&\colorbox{gray!20}{0.2238}&\colorbox{gray!20}{0.1948}&\colorbox{gray!20}{0.1310}&\colorbox{gray!20}{0.1801}&\colorbox{gray!20}{0.2033}&\colorbox{gray!20}{0.1382}&\colorbox{gray!20}{0.1871}&\colorbox{gray!20}{0.5317}&\colorbox{gray!20}{0.3732}&\colorbox{gray!20}{0.5419}&\colorbox{gray!20}{0.5280}&\colorbox{gray!20}{0.3659}&\colorbox{gray!20}{0.5380}



\\
                      & DBCNN \cite{quality:DBCNN}        & \colorbox{gray!20}{0.3227}&\colorbox{gray!20}{0.2200}&\colorbox{gray!20}{0.3453}&\colorbox{gray!20}{0.2421}&\colorbox{gray!20}{0.1591}&\colorbox{gray!20}{0.2323}&\colorbox{gray!20}{0.2218}&\colorbox{gray!20}{0.1480}&\colorbox{gray!20}{0.2486}&\colorbox{gray!20}{0.6688}&\colorbox{gray!20}{0.4743}&\colorbox{gray!20}{0.6579}&\colorbox{gray!20}{0.6285}&\colorbox{gray!20}{0.4306}&\colorbox{gray!20}{0.5990}



 \\
                      & QualiClip\cite{quality:qualiclip} &    \CLB{0.5795} & \CLB{0.4074} & \CLB{0.5388} & \CLB{0.4769} & \CLB{0.3258} & \CLB{0.4673} & \CLB{0.4846} & \CLB{0.3371} & \CLB{0.5112} & \CLB{0.7847} & \CLB{0.5849} & \CLB{0.7739} & \CLB{0.7691} & \CLB{0.5667} & \CLB{0.7420}



 \\
                      & TOPIQ-NR \cite{quality:topiq}     & \CLA{0.6443}&\CLA{0.4550}&\CLA{0.6295}&\CLA{0.6006}&\CLA{0.4248}&\CLA{0.5996}&\CLA{0.5703}&\CLA{0.4014}&\CLA{0.6153}&\CLA{0.8403}&\CLA{0.6454}&\CLA{0.8320}&\CLA{0.8322}&\CLA{0.6307}&\CLA{0.8241}


 \\ \bottomrule
\end{tabular}}

\resizebox{\linewidth}{!}{
\begin{tabular}{l|l|ccc|ccc|ccc|ccc|ccc}
\toprule
\multicolumn{2}{c|}{Dimension}           & \multicolumn{3}{c|}{Dis-level-1}  & \multicolumn{3}{c|}{Dis-level-2} & \multicolumn{3}{c|}{Dis-level-3}  & \multicolumn{3}{c|}{Dis-level-4}  & \multicolumn{3}{c}{Dis-level-5}  \\ \cline{1-17}
Group                 & \multicolumn{1}{c|}{Metric}        & SRCC$\uparrow$   & KRCC$\uparrow$   & PLCC$\uparrow$   & SRCC$\uparrow$   & KRCC$\uparrow$   & PLCC$\uparrow$   & SRCC$\uparrow$   & KRCC$\uparrow$   & PLCC$\uparrow$   & SRCC$\uparrow$   & KRCC$\uparrow$   & PLCC$\uparrow$   & SRCC$\uparrow$   & KRCC$\uparrow$   & PLCC$\uparrow$   \\ \midrule
\multirow{5}{*}{Zero} & PSNR          & 0.2405&0.1665&0.3445&0.2447&0.1644&0.2930&0.2261&0.1512&0.2695&0.2065&0.1357&0.2323&0.4268&0.2894&0.4473


 \\
                      & SSIM \cite{quality:ssim}         & 0.5085&0.3520&0.4753&0.4888&0.3389&0.4547&0.5601&0.3878&0.5299&0.5646&0.3852&0.5450&0.6977&0.4997&\textbf{0.6846}


 \\
                      & Brisque \cite{quality:brisque}      & 0.1471&0.0971&0.1276&0.2153&0.1466&0.2050&0.3106&0.2093&0.2767&0.3112&0.2058&0.2929&0.4560&0.3109&0.4301


 \\
                      & Q-Align \cite{quality:q-align}  & \textbf{0.6748}&\textbf{0.4887}&\textbf{0.6579}&\textbf{0.6826}&\textbf{0.4943}&\textbf{0.6539}&\textbf{0.7060}&\textbf{0.5079}&\textbf{0.6656}&\textbf{0.7316}&\textbf{0.5250}&\textbf{0.6937}&\textbf{0.7055}&\textbf{0.5040}&0.6775

 \\
                      & Q-Align+ \cite{quality:q-align+} & 0.4598&0.3136&0.4055&0.5730&0.4006&0.5013&0.5487&0.3758&0.4988&0.5392&0.3684&0.5110&0.4769&0.3238&0.4140


 \\ \cdashline{1-17}
\multirow{5}{*}{FR}   & AHIQ \cite{quality:ahiq}         & \CLB{0.7138} & \CLB{0.5226} & \CLB{0.7554} & \CLB{0.7917} & \CLB{0.5935} & \CLB{0.7843} & \CLB{0.7462} & \CLB{0.5506} & \CLB{0.7434} & \CLB{0.7828} & \CLB{0.5810} & \CLB{0.7761} & \CLA{0.8124} & \CLA{0.6122} & \CLB{0.8149}



 \\
                      & CKDN \cite{quality:ckdn}         & 0.6883&0.5044&0.7448&0.7236&0.5217&0.7186&\colorbox{gray!20}{0.7024}&0.5085&0.7073&\colorbox{gray!20}{0.6989}&\colorbox{gray!20}{0.5030}&\colorbox{gray!20}{0.6868}&\colorbox{gray!20}{0.6997}&0.5076&0.7045


 \\
                      & DISTS \cite{quality:dists}        & \colorbox{gray!20}{0.6557}&\colorbox{gray!20}{0.4749}&0.6641&\colorbox{gray!20}{0.6385}&\colorbox{gray!20}{0.4582}&\colorbox{gray!20}{0.6488}&\colorbox{gray!20}{0.6720}&\colorbox{gray!20}{0.4856}&0.6686&\colorbox{gray!20}{0.6875}&\colorbox{gray!20}{0.4945}&0.7040&0.7541&0.5540&0.7571


 \\
                      & LPIPS \cite{quality:lpips}        & \colorbox{gray!20}{0.6174}&\colorbox{gray!20}{0.4401}&\colorbox{gray!20}{0.6323}&\colorbox{gray!20}{0.6742}&\colorbox{gray!20}{0.4853}&\colorbox{gray!20}{0.6405}&\colorbox{gray!20}{0.6098}&\colorbox{gray!20}{0.4374}&\colorbox{gray!20}{0.5895}&\colorbox{gray!20}{0.6188}&\colorbox{gray!20}{0.4358}&\colorbox{gray!20}{0.5678}&\colorbox{gray!20}{0.6842}&\colorbox{gray!20}{0.4844}&\colorbox{gray!20}{0.6248}


 \\
                      & TOPIQ-FR \cite{quality:topiq}     & \CLA{0.7876}&\CLA{0.5949}&\CLA{0.8192}&\CLA{0.8030}&\CLA{0.6011}&\CLA{0.8001}&\CLA{0.7944}&\CLA{0.5941}&\CLA{0.7898}&\CLA{0.8161}&\CLA{0.6111}&\CLA{0.8051}&\CLB{0.8084}&\CLB{0.6033}&\CLA{0.8160}


 \\ \cdashline{1-17}
\multirow{6}{*}{NR}   & CLIPIQA \cite{quality:CLIPIQA}       & \colorbox{gray!20}{0.2597}&\colorbox{gray!20}{0.1727}&\colorbox{gray!20}{0.2715}&\colorbox{gray!20}{0.3258}&\colorbox{gray!20}{0.2201}&\colorbox{gray!20}{0.3295}&\colorbox{gray!20}{0.2608}&\colorbox{gray!20}{0.1768}&\colorbox{gray!20}{0.2473}&\colorbox{gray!20}{0.3331}&\colorbox{gray!20}{0.2252}&\colorbox{gray!20}{0.3480}&\colorbox{gray!20}{0.4283}&\colorbox{gray!20}{0.2921}&\colorbox{gray!20}{0.4228}



 \\
                      & CNNIQA \cite{quality:CNNIQA}       & \colorbox{gray!20}{0.3490}&\colorbox{gray!20}{0.2356}&\colorbox{gray!20}{0.3572}&\colorbox{gray!20}{0.3909}&\colorbox{gray!20}{0.2676}&\colorbox{gray!20}{0.3737}&\colorbox{gray!20}{0.5042}&\colorbox{gray!20}{0.3497}&\colorbox{gray!20}{0.4942}&\colorbox{gray!20}{0.5044}&\colorbox{gray!20}{0.3412}&\colorbox{gray!20}{0.5092}&\colorbox{gray!20}{0.6176}&\colorbox{gray!20}{0.4311}&\colorbox{gray!20}{0.6028}



 \\
                      & DBCNN \cite{quality:DBCNN}        & \colorbox{gray!20}{0.5077}&\colorbox{gray!20}{0.3527}&\colorbox{gray!20}{0.5513}&\colorbox{gray!20}{0.5553}&\colorbox{gray!20}{0.3855}&\colorbox{gray!20}{0.5580}&\colorbox{gray!20}{0.6177}&\colorbox{gray!20}{0.4259}&\colorbox{gray!20}{0.6008}&\colorbox{gray!20}{0.6373}&\colorbox{gray!20}{0.4404}&\colorbox{gray!20}{0.6217}&\colorbox{gray!20}{0.6472}&\colorbox{gray!20}{0.4496}&\colorbox{gray!20}{0.6435}



\\
                      & QualiClip\cite{quality:qualiclip}   &\CLB{0.7072} & \CLB{0.5115} & \CLB{0.7188} & \CLB{0.7337} & \CLB{0.5366} & \CLB{0.7187} & \colorbox{gray!20}{\CLB{0.7027}} & \colorbox{gray!20}{\CLB{0.5048}} & \CLB{0.6866} & \CLB{0.7602} & \CLB{0.5537} & \CLB{0.7493} & \CLB{0.7352} & \CLB{0.5418} & \CLB{0.7379}



 \\
                      & TOPIQ-NR \cite{quality:topiq}     & \CLA{0.7788}&\CLA{0.5818}&\CLA{0.8044}&\CLA{0.7974}&\CLA{0.6014}&\CLA{0.7993}&\CLA{0.7966}&\CLA{0.6011}&\CLA{0.7911}&\CLA{0.8071}&\CLA{0.6035}&\CLA{0.7967}&\CLA{0.7965}&\CLA{0.5952}&\CLA{0.8027}


 \\ \bottomrule
\end{tabular}}
\vspace{-4mm}
\end{table*}

\section{Low-level Attribute Distribution}
% Jiahao，周一说

\begin{figure}[t]
    \centering
    \includegraphics[width=\linewidth]{supp/feature2.pdf}
    \caption{Low-level feature distribution of MPD, normalized and visualized in 30 corruption subsets. Different colors denote \GroupA{Luminance}, \GroupB{Contrast}, \GroupC{Chrominance}, \GroupD{Blur}, and \GroupE{Spatial Information}.}
    \label{fig:low}
\end{figure}

Figure \ref{fig:low} shows the distribution of low-level features of all instances of Embodied IQA. After overall regularization, 30 types of corruption are grouped and displayed. The features considered include Luminance, Contrast, Chrominance, Blur, and Spatial Information. There are significant differences in these low-level attributes for different corruptions. For example, in the first three blurry cases, the blur curve is left-biased and becomes right-biased after sharpening. In general, similar corruption categories will lead to similar results (such as five noise-related and four block-related). Two of the denoises have the sharpest distributions; Color quantization, Grayscale quantization, Sharpness change, and Contrast change are the most irregular. These findings deserve further exploration.

\begin{figure*}[t]
    \centering
    \begin{minipage}[]{0.49\linewidth}
    \centering
    \centerline{\includegraphics[width = \textwidth]{supp/case1.pdf}}
    % \centerline{VQA}\medskip
    \end{minipage}
    \begin{minipage}[]{0.49\linewidth}
    \centering
    \centerline{\includegraphics[width = \textwidth]{supp/case2.pdf}}
    % \centerline{CAP}\medskip
    \end{minipage}

    \begin{minipage}[]{0.49\linewidth}
    \centering
    \centerline{\includegraphics[width = \textwidth]{supp/case3.pdf}}
    % \centerline{CAP}\medskip
    \end{minipage}
    \begin{minipage}[]{0.49\linewidth}
    \centering
    \centerline{\includegraphics[width = \textwidth]{supp/case4.pdf}}
    % \centerline{CAP}\medskip
    \end{minipage}

    \caption{Positive and negative cases. Slight distortion may significantly affect the inference result of Embodied AI, while severe distortion may not. Emphasizing the significance of Embodied IQA.}
    \label{fig:case}
\end{figure*}

\begin{figure}[t]
    \centering
    \includegraphics[width=\linewidth]{supp/corruption_down.pdf}
    \caption{Visualization of 30 distortion types. Strength from left (Level 1) to right (Level 5).}
    \label{fig:distort}
    \vspace{-2mm}
\end{figure}

\section{Cases Study}
% Jiahao (图，task1啥，vlm输出啥,vla输出啥) (高分低分样例，VLM分和VLA分)
Figure \ref{fig:case} shows four typical examples from the Embodied IQA database (center-cropped for visualization), including the VLA inference results for reference/distorted image pairs under different distortions. The results in the upper left and lower right corners are as expected, the more severe the distortion, the lower the score. However, for the distortion level 5 in the upper righ, since it does not affect any objects on the desktop, the `gray block' as the target of the task is not affected, so the subjective score is as high as 4.52; on the contrary, although the distortion level in the lower left corner is only 1, the `lost macro block' happens to be the target object, so the VLA Position and Rotation are greatly changed with a score only 2.98.
Figure \ref{fig:distort} shows 30 distortion types at different strength levels from 1 to 5. In the previous human-oriented scenario, the visual quality of different corruptions is similar at the same strength. However, from the example above, the preference of Embodied AI depends on the task, which significantly differs from traditional IQA paradigm.  We hope that our database can further inspire better quality metrics for Embodied AI.

\section{Disclaimer}
% Chunyi
The main purpose of this study is to apply IQA to Embodied AI to promote its Real-world application, rather than to praise or criticize any VLM, VLA, or IQA model.
We evaluate image samples rather than models. Lower scores do not mean that the performance of downstream VLM/VLA is poor, but distortion has a greater impact on it; similarly, lower correlation coefficients do not mean defects in the IQA method, but rather indicate the huge difference between Embodied and traditional IQA.
Considering the scale of the database, we will open it in several stages for non-commercial use, and sincerely hope that future robotic-oriented IQA metrics can drive the development of Embodied AI.



%%%%%%%%%%%%%%%%%%%%%%%%%%%%%%%%%%%%%%%%%%%%%%%%%%%%%%%%%%%%

% \newpage
% \section*{NeurIPS Paper Checklist}

% %%% BEGIN INSTRUCTIONS %%%
% % The checklist is designed to encourage best practices for responsible machine learning research, addressing issues of reproducibility, transparency, research ethics, and societal impact. Do not remove the checklist: {\bf The papers not including the checklist will be desk rejected.} The checklist should follow the references and follow the (optional) supplemental material.  The checklist does NOT count towards the page
% % limit. 

% % Please read the checklist guidelines carefully for information on how to answer these questions. For each question in the checklist:
% % \begin{itemize}
% %     \item You should answer \answerYes{}, \answerNo{}, or \answerNA{}.
% %     \item \answerNA{} means either that the question is Not Applicable for that particular paper or the relevant information is Not Available.
% %     \item Please provide a short (1–2 sentence) justification right after your answer (even for NA). 
% %    % \item {\bf The papers not including the checklist will be desk rejected.}
% % \end{itemize}

% % {\bf The checklist answers are an integral part of your paper submission.} They are visible to the reviewers, area chairs, senior area chairs, and ethics reviewers. You will be asked to also include it (after eventual revisions) with the final version of your paper, and its final version will be published with the paper.

% % The reviewers of your paper will be asked to use the checklist as one of the factors in their evaluation. While "\answerYes{}" is generally preferable to "\answerNo{}", it is perfectly acceptable to answer "\answerNo{}" provided a proper justification is given (e.g., "error bars are not reported because it would be too computationally expensive" or "we were unable to find the license for the dataset we used"). In general, answering "\answerNo{}" or "\answerNA{}" is not grounds for rejection. While the questions are phrased in a binary way, we acknowledge that the true answer is often more nuanced, so please just use your best judgment and write a justification to elaborate. All supporting evidence can appear either in the main paper or the supplemental material, provided in appendix. If you answer \answerYes{} to a question, in the justification please point to the section(s) where related material for the question can be found.

% % IMPORTANT, please:
% % \begin{itemize}
% %     \item {\bf Delete this instruction block, but keep the section heading ``NeurIPS Paper Checklist"},
% %     \item  {\bf Keep the checklist subsection headings, questions/answers and guidelines below.}
% %     \item {\bf Do not modify the questions and only use the provided macros for your answers}.
% % \end{itemize} 
 

% %%% END INSTRUCTIONS %%%


% \begin{enumerate}

% \item {\bf Claims}
%     \item[] Question: Do the main claims made in the abstract and introduction accurately reflect the paper's contributions and scope?
%     \item[] Answer: \answerYes{}
%     \item[] Justification: Abstract summarizes our three main contributions, including theory, data, and experiments. The specific content is highlighted in the list in the introduction.
%     \item[] Guidelines:
%     \begin{itemize}
%         \item The answer NA means that the abstract and introduction do not include the claims made in the paper.
%         \item The abstract and/or introduction should clearly state the claims made, including the contributions made in the paper and important assumptions and limitations. A No or NA answer to this question will not be perceived well by the reviewers. 
%         \item The claims made should match theoretical and experimental results, and reflect how much the results can be expected to generalize to other settings. 
%         \item It is fine to include aspirational goals as motivation as long as it is clear that these goals are not attained by the paper. 
%     \end{itemize}

% \item {\bf Limitations}
%     \item[] Question: Does the paper discuss the limitations of the work performed by the authors?
%     \item[] Answer: \answerYes{}
%     \item[] Justification: As the first Embodied IQA work, considering the high cost of Real-world experiments, we limit perception to vision (accounting for 80\% of the total signal input) and execution to upper limbs (accounting for 50\% of the total motor action). We look forward to further improvement in subsequent work, see Supplementary Section A for details.
%     \item[] Guidelines:
%     \begin{itemize}
%         \item The answer NA means that the paper has no limitation while the answer No means that the paper has limitations, but those are not discussed in the paper. 
%         \item The authors are encouraged to create a separate "Limitations" section in their paper.
%         \item The paper should point out any strong assumptions and how robust the results are to violations of these assumptions (e.g., independence assumptions, noiseless settings, model well-specification, asymptotic approximations only holding locally). The authors should reflect on how these assumptions might be violated in practice and what the implications would be.
%         \item The authors should reflect on the scope of the claims made, e.g., if the approach was only tested on a few datasets or with a few runs. In general, empirical results often depend on implicit assumptions, which should be articulated.
%         \item The authors should reflect on the factors that influence the performance of the approach. For example, a facial recognition algorithm may perform poorly when image resolution is low or images are taken in low lighting. Or a speech-to-text system might not be used reliably to provide closed captions for online lectures because it fails to handle technical jargon.
%         \item The authors should discuss the computational efficiency of the proposed algorithms and how they scale with dataset size.
%         \item If applicable, the authors should discuss possible limitations of their approach to address problems of privacy and fairness.
%         \item While the authors might fear that complete honesty about limitations might be used by reviewers as grounds for rejection, a worse outcome might be that reviewers discover limitations that aren't acknowledged in the paper. The authors should use their best judgment and recognize that individual actions in favor of transparency play an important role in developing norms that preserve the integrity of the community. Reviewers will be specifically instructed to not penalize honesty concerning limitations.
%     \end{itemize}

% \item {\bf Theory assumptions and proofs}
%     \item[] Question: For each theoretical result, does the paper provide the full set of assumptions and a complete (and correct) proof?
%     \item[] Answer: \answerYes{}
%     \item[] Justification: The findings of this paper are based on the assumption of the Mertonian system, namelythe Perception-Cognition-Decision-Execution pipeline of Embodied AI. This assumption has been widely accepted in the field of robotics (see Section 2.1), and Figures 4-8 all imply the differences between these four steps, proving that these steps cannot be viewed as a whole.
%     \item[] Guidelines:
%     \begin{itemize}
%         \item The answer NA means that the paper does not include theoretical results. 
%         \item All the theorems, formulas, and proofs in the paper should be numbered and cross-referenced.
%         \item All assumptions should be clearly stated or referenced in the statement of any theorems.
%         \item The proofs can either appear in the main paper or the supplemental material, but if they appear in the supplemental material, the authors are encouraged to provide a short proof sketch to provide intuition. 
%         \item Inversely, any informal proof provided in the core of the paper should be complemented by formal proofs provided in appendix or supplemental material.
%         \item Theorems and Lemmas that the proof relies upon should be properly referenced. 
%     \end{itemize}

%     \item {\bf Experimental result reproducibility}
%     \item[] Question: Does the paper fully disclose all the information needed to reproduce the main experimental results of the paper to the extent that it affects the main claims and/or conclusions of the paper (regardless of whether the code and data are provided or not)?
%     \item[] Answer: \answerYes{}
%     \item[] Justification: Section 5.1 lists the details of reproducing the IQA experiment. The results in Section 5.2 are consistent with the main claims and conclusions.
%     \item[] Guidelines:
%     \begin{itemize}
%         \item The answer NA means that the paper does not include experiments.
%         \item If the paper includes experiments, a No answer to this question will not be perceived well by the reviewers: Making the paper reproducible is important, regardless of whether the code and data are provided or not.
%         \item If the contribution is a dataset and/or model, the authors should describe the steps taken to make their results reproducible or verifiable. 
%         \item Depending on the contribution, reproducibility can be accomplished in various ways. For example, if the contribution is a novel architecture, describing the architecture fully might suffice, or if the contribution is a specific model and empirical evaluation, it may be necessary to either make it possible for others to replicate the model with the same dataset, or provide access to the model. In general. releasing code and data is often one good way to accomplish this, but reproducibility can also be provided via detailed instructions for how to replicate the results, access to a hosted model (e.g., in the case of a large language model), releasing of a model checkpoint, or other means that are appropriate to the research performed.
%         \item While NeurIPS does not require releasing code, the conference does require all submissions to provide some reasonable avenue for reproducibility, which may depend on the nature of the contribution. For example
%         \begin{enumerate}
%             \item If the contribution is primarily a new algorithm, the paper should make it clear how to reproduce that algorithm.
%             \item If the contribution is primarily a new model architecture, the paper should describe the architecture clearly and fully.
%             \item If the contribution is a new model (e.g., a large language model), then there should either be a way to access this model for reproducing the results or a way to reproduce the model (e.g., with an open-source dataset or instructions for how to construct the dataset).
%             \item We recognize that reproducibility may be tricky in some cases, in which case authors are welcome to describe the particular way they provide for reproducibility. In the case of closed-source models, it may be that access to the model is limited in some way (e.g., to registered users), but it should be possible for other researchers to have some path to reproducing or verifying the results.
%         \end{enumerate}
%     \end{itemize}


% \item {\bf Open access to data and code}
%     \item[] Question: Does the paper provide open access to the data and code, with sufficient instructions to faithfully reproduce the main experimental results, as described in supplemental material?
%     \item[] Answer: \answerYes{}
%     \item[] Justification: All data has been prepared and will be open-sourced, including the VLM/VLA score, the IQA configuration file, and the reference/distorted image collection toolkit.
%     \item[] Guidelines:
%     \begin{itemize}
%         \item The answer NA means that paper does not include experiments requiring code.
%         \item Please see the NeurIPS code and data submission guidelines (\url{https://nips.cc/public/guides/CodeSubmissionPolicy}) for more details.
%         \item While we encourage the release of code and data, we understand that this might not be possible, so “No” is an acceptable answer. Papers cannot be rejected simply for not including code, unless this is central to the contribution (e.g., for a new open-source benchmark).
%         \item The instructions should contain the exact command and environment needed to run to reproduce the results. See the NeurIPS code and data submission guidelines (\url{https://nips.cc/public/guides/CodeSubmissionPolicy}) for more details.
%         \item The authors should provide instructions on data access and preparation, including how to access the raw data, preprocessed data, intermediate data, and generated data, etc.
%         \item The authors should provide scripts to reproduce all experimental results for the new proposed method and baselines. If only a subset of experiments are reproducible, they should state which ones are omitted from the script and why.
%         \item At submission time, to preserve anonymity, the authors should release anonymized versions (if applicable).
%         \item Providing as much information as possible in supplemental material (appended to the paper) is recommended, but including URLs to data and code is permitted.
%     \end{itemize}


% \item {\bf Experimental setting/details}
%     \item[] Question: Does the paper specify all the training and test details (e.g., data splits, hyperparameters, how they were chosen, type of optimizer, etc.) necessary to understand the results?
%     \item[] Answer: \answerYes{}
%     \item[] Justification: Section 5.1 includes training and test details in simulation, see Supplementary Section B for further Real-world details.
%     \item[] Guidelines:
%     \begin{itemize}
%         \item The answer NA means that the paper does not include experiments.
%         \item The experimental setting should be presented in the core of the paper to a level of detail that is necessary to appreciate the results and make sense of them.
%         \item The full details can be provided either with the code, in appendix, or as supplemental material.
%     \end{itemize}

% \item {\bf Experiment statistical significance}
%     \item[] Question: Does the paper report error bars suitably and correctly defined or other appropriate information about the statistical significance of the experiments?
%     \item[] Answer: \answerYes{}
%     \item[] Justification: The error bars (confidence level 0.9) are included in the Embodied score distribution, see Figure 6.
%     \item[] Guidelines:
%     \begin{itemize}
%         \item The answer NA means that the paper does not include experiments.
%         \item The authors should answer "Yes" if the results are accompanied by error bars, confidence intervals, or statistical significance tests, at least for the experiments that support the main claims of the paper.
%         \item The factors of variability that the error bars are capturing should be clearly stated (for example, train/test split, initialization, random drawing of some parameter, or overall run with given experimental conditions).
%         \item The method for calculating the error bars should be explained (closed form formula, call to a library function, bootstrap, etc.)
%         \item The assumptions made should be given (e.g., Normally distributed errors).
%         \item It should be clear whether the error bar is the standard deviation or the standard error of the mean.
%         \item It is OK to report 1-sigma error bars, but one should state it. The authors should preferably report a 2-sigma error bar than state that they have a 96\% CI, if the hypothesis of Normality of errors is not verified.
%         \item For asymmetric distributions, the authors should be careful not to show in tables or figures symmetric error bars that would yield results that are out of range (e.g. negative error rates).
%         \item If error bars are reported in tables or plots, The authors should explain in the text how they were calculated and reference the corresponding figures or tables in the text.
%     \end{itemize}

% \item {\bf Experiments compute resources}
%     \item[] Question: For each experiment, does the paper provide sufficient information on the computer resources (type of compute workers, memory, time of execution) needed to reproduce the experiments?
%     \item[] Answer: \answerYes{}
%     \item[] Justification: Section 5.1 includes the specific version of the robotic arm and computational resources, see Supplementary Section B for detailed parameters.
%     \item[] Guidelines:
%     \begin{itemize}
%         \item The answer NA means that the paper does not include experiments.
%         \item The paper should indicate the type of compute workers CPU or GPU, internal cluster, or cloud provider, including relevant memory and storage.
%         \item The paper should provide the amount of compute required for each of the individual experimental runs as well as estimate the total compute. 
%         \item The paper should disclose whether the full research project required more compute than the experiments reported in the paper (e.g., preliminary or failed experiments that didn't make it into the paper). 
%     \end{itemize}
    
% \item {\bf Code of ethics}
%     \item[] Question: Does the research conducted in the paper conform, in every respect, with the NeurIPS Code of Ethics \url{https://neurips.cc/public/EthicsGuidelines}?
%     \item[] Answer: \answerYes{}
%     \item[] Justification:  This study complies with the ethical standards of the NeurIPS Code of Ethics and Helsinki Declaration. 
%     \item[] Guidelines:
%     \begin{itemize}
%         \item The answer NA means that the authors have not reviewed the NeurIPS Code of Ethics.
%         \item If the authors answer No, they should explain the special circumstances that require a deviation from the Code of Ethics.
%         \item The authors should make sure to preserve anonymity (e.g., if there is a special consideration due to laws or regulations in their jurisdiction).
%     \end{itemize}


% \item {\bf Broader impacts}
%     \item[] Question: Does the paper discuss both potential positive societal impacts and negative societal impacts of the work performed?
%     \item[] Answer: \answerYes{}
%     \item[] Justification: The positive impact is to expand the real-world applications of embedded AI. However, considering the gap between human and robot visual systems, the widespread use of this quality indicator may affect the original human-oriented tasks. See Supplementary Section A for details.
%     \item[] Guidelines:
%     \begin{itemize}
%         \item The answer NA means that there is no societal impact of the work performed.
%         \item If the authors answer NA or No, they should explain why their work has no societal impact or why the paper does not address societal impact.
%         \item Examples of negative societal impacts include potential malicious or unintended uses (e.g., disinformation, generating fake profiles, surveillance), fairness considerations (e.g., deployment of technologies that could make decisions that unfairly impact specific groups), privacy considerations, and security considerations.
%         \item The conference expects that many papers will be foundational research and not tied to particular applications, let alone deployments. However, if there is a direct path to any negative applications, the authors should point it out. For example, it is legitimate to point out that an improvement in the quality of generative models could be used to generate deepfakes for disinformation. On the other hand, it is not needed to point out that a generic algorithm for optimizing neural networks could enable people to train models that generate Deepfakes faster.
%         \item The authors should consider possible harms that could arise when the technology is being used as intended and functioning correctly, harms that could arise when the technology is being used as intended but gives incorrect results, and harms following from (intentional or unintentional) misuse of the technology.
%         \item If there are negative societal impacts, the authors could also discuss possible mitigation strategies (e.g., gated release of models, providing defenses in addition to attacks, mechanisms for monitoring misuse, mechanisms to monitor how a system learns from feedback over time, improving the efficiency and accessibility of ML).
%     \end{itemize}
    
% \item {\bf Safeguards}
%     \item[] Question: Does the paper describe safeguards that have been put in place for responsible release of data or models that have a high risk for misuse (e.g., pretrained language models, image generators, or scraped datasets)?
%     \item[] Answer: \answerNA{}
%     \item[] Justification: The paper poses no such risks.
%     \item[] Guidelines:
%     \begin{itemize}
%         \item The answer NA means that the paper poses no such risks.
%         \item Released models that have a high risk for misuse or dual-use should be released with necessary safeguards to allow for controlled use of the model, for example by requiring that users adhere to usage guidelines or restrictions to access the model or implementing safety filters. 
%         \item Datasets that have been scraped from the Internet could pose safety risks. The authors should describe how they avoided releasing unsafe images.
%         \item We recognize that providing effective safeguards is challenging, and many papers do not require this, but we encourage authors to take this into account and make a best faith effort.
%     \end{itemize}

% \item {\bf Licenses for existing assets}
%     \item[] Question: Are the creators or original owners of assets (e.g., code, data, models), used in the paper, properly credited and are the license and terms of use explicitly mentioned and properly respected?
%     \item[] Answer: \answerYes{}
%     \item[] Justification: All the original data are either self-captured, or allowed to be rearranged/adapted/re-created for non-commercial purposes.
%     \item[] Guidelines:
%     \begin{itemize}
%         \item The answer NA means that the paper does not use existing assets.
%         \item The authors should cite the original paper that produced the code package or dataset.
%         \item The authors should state which version of the asset is used and, if possible, include a URL.
%         \item The name of the license (e.g., CC-BY 4.0) should be included for each asset.
%         \item For scraped data from a particular source (e.g., website), the copyright and terms of service of that source should be provided.
%         \item If assets are released, the license, copyright information, and terms of use in the package should be provided. For popular datasets, \url{paperswithcode.com/datasets} has curated licenses for some datasets. Their licensing guide can help determine the license of a dataset.
%         \item For existing datasets that are re-packaged, both the original license and the license of the derived asset (if it has changed) should be provided.
%         \item If this information is not available online, the authors are encouraged to reach out to the asset's creators.
%     \end{itemize}

% \item {\bf New assets}
%     \item[] Question: Are new assets introduced in the paper well documented and is the documentation provided alongside the assets?
%     \item[] Answer: \answerNA{}
%     \item[] Justification: All data are under CC BY-NC-SA 4.0  license for non-commercial use.
%     \item[] Guidelines:
%     \begin{itemize}
%         \item The answer NA means that the paper does not release new assets.
%         \item Researchers should communicate the details of the dataset/code/model as part of their submissions via structured templates. This includes details about training, license, limitations, etc. 
%         \item The paper should discuss whether and how consent was obtained from people whose asset is used.
%         \item At submission time, remember to anonymize your assets (if applicable). You can either create an anonymized URL or include an anonymized zip file.
%     \end{itemize}

% \item {\bf Crowdsourcing and research with human subjects}
%     \item[] Question: For crowdsourcing experiments and research with human subjects, does the paper include the full text of instructions given to participants and screenshots, if applicable, as well as details about compensation (if any)? 
%     \item[] Answer: \answerYes{}
%     \item[] Justification: The full text of instructions given to participants and screenshots ares included, see Supplementary Section C for detail.
%     \item[] Guidelines:
%     \begin{itemize}
%         \item The answer NA means that the paper does not involve crowdsourcing nor research with human subjects.
%         \item Including this information in the supplemental material is fine, but if the main contribution of the paper involves human subjects, then as much detail as possible should be included in the main paper. 
%         \item According to the NeurIPS Code of Ethics, workers involved in data collection, curation, or other labor should be paid at least the minimum wage in the country of the data collector. 
%     \end{itemize}

% \item {\bf Institutional review board (IRB) approvals or equivalent for research with human subjects}
%     \item[] Question: Does the paper describe potential risks incurred by study participants, whether such risks were disclosed to the subjects, and whether Institutional Review Board (IRB) approvals (or an equivalent approval/review based on the requirements of your country or institution) were obtained?
%     \item[] Answer: \answerYes{}
%     \item[] Justification: This study protects the rights and welfare of the participants involved, with appropriate measures to maintain the confidentiality and security of participant data, see Supplementary Section C for detail.
%     \item[] Guidelines:
%     \begin{itemize}
%         \item The answer NA means that the paper does not involve crowdsourcing nor research with human subjects.
%         \item Depending on the country in which research is conducted, IRB approval (or equivalent) may be required for any human subjects research. If you obtained IRB approval, you should clearly state this in the paper. 
%         \item We recognize that the procedures for this may vary significantly between institutions and locations, and we expect authors to adhere to the NeurIPS Code of Ethics and the guidelines for their institution. 
%         \item For initial submissions, do not include any information that would break anonymity (if applicable), such as the institution conducting the review.
%     \end{itemize}

% \item {\bf Declaration of LLM usage}
%     \item[] Question: Does the paper describe the usage of LLMs if it is an important, original, or non-standard component of the core methods in this research? Note that if the LLM is used only for writing, editing, or formatting purposes and does not impact the core methodology, scientific rigorousness, or originality of the research, declaration is not required.
%     %this research? 
%     \item[] Answer: \answerYes{}
%     \item[] Justification: LLM is only used for grammar check.
%     \item[] Guidelines:
%     \begin{itemize}
%         \item The answer NA means that the core method development in this research does not involve LLMs as any important, original, or non-standard components.
%         \item Please refer to our LLM policy (\url{https://neurips.cc/Conferences/2025/LLM}) for what should or should not be described.
%     \end{itemize}

% \end{enumerate}

{\small
\bibliographystyle{unsrt}
\bibliography{neurips_2025}
}

\end{document}